\documentclass[a4paper,11pt]{article}

% --- Pacchetti Fondamentali ---
\usepackage[utf8]{inputenc}
\usepackage[T1]{fontenc}
\usepackage[italian]{babel}
\usepackage{geometry}
\geometry{top=2.5cm, bottom=2.5cm, left=2.5cm, right=2.5cm}

% --- Font Moderni ---
\usepackage{lmodern}
\usepackage{helvet}
\renewcommand{\familydefault}{\sfdefault} % Usa sans-serif per tutto il documento

% --- Colori e Grafica ---
\usepackage{xcolor}
\usepackage{tikz}
\usepackage{amssymb}
\usetikzlibrary{shapes, arrows.meta, positioning, chains, fit, calc, matrix, trees, backgrounds, shadows}

% --- Definizione Colori (Light Theme + Orange + New Cover Colors)
\definecolor{mainOrange}{RGB}{255, 140, 0}
\definecolor{jsYellow}{HTML}{F7DF1E}
\definecolor{codeBg}{gray}{0.95}
\definecolor{codeBlue}{HTML}{0000FF}
\definecolor{codePink}{HTML}{800080}
\definecolor{codeGreen}{HTML}{008000}
\definecolor{niceBlue}{HTML}{2E86C1}
\definecolor{jsDark}{HTML}{282C34}    % Grigio scuro (stile editor codice)
\definecolor{tsBlue}{HTML}{3178C6}    % Blu TypeScript
\definecolor{textGray}{RGB}{100, 100, 100}

% --- PALETTE COLORI ACCADEMICA (Nuovo Capitolo) ---
\definecolor{paletteIndigo}{HTML}{4B0082}
\definecolor{paletteRegalia}{HTML}{522D80}
\definecolor{paletteLenurple}{HTML}{BA55D3}
\definecolor{paletteLavender}{HTML}{E6E6FA}

% --- COMANDO COPERTINA JS (Pure TikZ - No Geometry Hacks) ---
\newcommand{\jscover}[1]{
    \begin{titlepage}
    \thispagestyle{empty}
    \begin{tikzpicture}[remember picture, overlay]
        
        % 1. Sfondo Laterale (Scuro)
        \fill[jsDark] (current page.north west) rectangle ($(current page.south west) + (4cm,0)$);
        
        % 2. Linea Accento (Giallo)
        \draw[jsYellow, line width=4pt] ($(current page.north west) + (4.2cm,-2cm)$) -- ($(current page.south west) + (4.2cm,2cm)$);
        
        % 3. Contenuto Testuale (Posizionato assolutamente con TikZ Node)
        % CENTRATURA MATEMATICA:
        % Larghezza Pagina: 21cm
        % Sidebar: 4cm
        % Spazio Bianco Rimanente: 17cm
        % Centro dello Spazio Bianco: 4cm + (17cm / 2) = 12.5cm
        % NOTE: Riduciamo text width a 12cm per margini più ampi (2.5cm per lato) e visuale migliore
        \node[anchor=north, inner sep=0, text width=12cm, align=center] at ($(current page.north west) + (12.5cm,-5cm)$) {
            
                {\fontsize{18}{22}\selectfont \textcolor{textGray}{Università di Pisa}}\\[0.2cm]
                {\fontsize{16}{20}\selectfont \textcolor{textGray}{Corso di Laurea in Informatica}}\\[2cm]
                
                {\fontsize{60}{70}\selectfont \textbf{\textcolor{jsDark}{LA BIBBIA}}}\\[0.5cm]
                
                {\fontsize{35}{45}\selectfont \textbf{\textcolor{black}{DI LABORATORIO}}}\\[0.2cm]
                {\fontsize{35}{45}\selectfont \textbf{\textcolor{black}{JAVASCRIPT}}}\\[1.5cm]
                
                % Linea divisoria Gialla (Centrata)
                \textcolor{jsYellow}{\rule{10cm}{4pt}}\\[1cm]
                
                {\Large \textbf{Algoritmi, Strutture Dati e TypeScript}}\\[0.2cm]
                {\large \textit{Stack, Queue, Alberi, Ricorsione e Higher-Order Functions}}\\[1.5cm]
                
                {\large Docenti titolari:}\\[0.2cm]
                {\Large \textbf{#1}}\\[2cm]
                
                \textbf{\large Autore del riassunto:}\\
                {\large Joseph Zucchelli (IlBroccolo)}\\[2cm]
                
                \vfill
                {\large \textbf{Anno Accademico 2025 - 2026}}
            
        };
        
        % Badge JS
        \node[anchor=south east] at ($(current page.south east) + (-2cm, 2cm)$) {
            \tikz\draw[fill=jsYellow, draw=none] (0,0) rectangle (2,2) node[midway, black, font=\bfseries\huge] {JS};
        };

    \end{tikzpicture}
    \end{titlepage}
    \clearpage
}

% --- Box per il Codice e Note (tcolorbox) ---
\usepackage[most]{tcolorbox}
\usepackage{listings}
\usepackage{titlesec} % Per colorare i titoli

% Setup Titoli Arancioni
\titleformat{\section}
{\color{mainOrange}\normalfont\Large\bfseries}
{\color{mainOrange}\thesection}{1em}{}[\color{mainOrange}\hrule height 1pt]

\titleformat{\subsection}
{\color{mainOrange}\normalfont\large\bfseries}
{\color{mainOrange}\thesubsection}{1em}{}

% Stile per il codice JS
\lstdefinelanguage{JavaScript}{
  keywords={typeof, new, true, false, catch, function, return, null, catch, switch, var, if, in, while, do, else, case, break, const, let, class, extends, super, static, get, set, yield, constructor, instanceof},
  keywordstyle=\color{blue}\bfseries,
  ndkeywords={export, boolean, throw, implements, import, this, console, log, prompt, alert},
  ndkeywordstyle=\color{teal}\bfseries,
  identifierstyle=\color{black}, 
  sensitive=false,
  comment=[l]{//},
  morecomment=[s]{/*}{*/},
  commentstyle=\color{green!50!black}\ttfamily,
  stringstyle=\color{red!60!black}\ttfamily,
  morestring=[b]',
  morestring=[b]",
  literate={à}{{\`a}}1
           {è}{{\`e}}1
           {é}{{\'e}}1
           {ì}{{\`i}}1
           {ò}{{\`o}}1
           {ù}{{\`u}}1
}

% Box Codice
\newtcblisting{jsbox}[1][]{
  colback=codeBg,
  colframe=gray!50,
  listing only,
  listing options={
    language=JavaScript,
    basicstyle=\ttfamily\small, % Testo nero di default
    numbers=left,
    numberstyle=\tiny\color{gray},
    breaklines=true,
    columns=fullflexible
  },
  title={\textbf{JS Code}},
  coltitle=black,
  fonttitle=\bfseries\ttfamily,
  enhanced,
  attach boxed title to top right={yshift=-3mm, xshift=-3mm},
  boxed title style={colback=jsYellow},
  #1
}

% Box Funzioni Viola (Richiesta Utente)
\newtcblisting{funcbox}[1][]{
  colback=violet!5,        % Sfondo Viola chiarissimo
  colframe=violet!75,      % Bordo Viola
  listing only,
  listing options={
    language=JavaScript,
    basicstyle=\ttfamily\small,
    numbers=left,
    numberstyle=\tiny\color{gray},
    breaklines=true,
    columns=fullflexible
  },
  title={\textbf{JS Function: #1}},
  coltitle=white,
  fonttitle=\bfseries\ttfamily,
  enhanced,
  attach boxed title to top left={yshift=-3mm, xshift=3mm},
  boxed title style={colback=violet}, 
  drop shadow,
  #1
}

% Box Concetti/Note
\newtcolorbox{concept}[1]{
  colback=niceBlue!5,
  colframe=niceBlue,
  title={#1},
  fonttitle=\bfseries,
  sharp corners=south,
  enhanced,
  drop shadow
}

% --- NUOVI STILI BOX STANDARDIZZATI ---
\tcbuselibrary{skins, breakable}

\newtcolorbox{definition}[1]{
    colback=mainOrange!5,
    colframe=mainOrange,
    fonttitle=\bfseries,
    title={Definizione: #1},
    sharp corners=south,
    enhanced,
    breakable,
    drop shadow
}

\newtcolorbox{algo}[1]{
    colback=yellow!10,
    colframe=orange!80!black,
    fonttitle=\bfseries,
    title={Algoritmo: #1},
    enhanced,
    breakable,
    drop shadow
}

\newtcolorbox{curiosita}[1]{
    colback=violet!5,
    colframe=violet,
    fonttitle=\bfseries,
    title={Nota a margine: #1},
    enhanced,
    breakable,
    drop shadow
}

\newtcolorbox{warning}[1]{
    colback=red!5,
    colframe=red!60!black,
    fonttitle=\bfseries,
    title={Attenzione: #1},
    enhanced,
    breakable,
    drop shadow
}

% --- LINK E PDF SETUP ---
\usepackage{hyperref}
\hypersetup{
    colorlinks=true,
    linkcolor=black,
    filecolor=magenta,      
    urlcolor=cyan,
    pdftitle={Bibbia JavaScript},
    pdfpagemode=FullScreen,
}

\begin{document}
\pagenumbering{gobble} % Niente numeri sulla copertina

\jscover{Prof. Jacopo Soldani \\ Prof. Davide Bacciu} % Nomi corretti
\clearpage
\pagenumbering{arabic} % Numeri arabi dall'indice in poi
\tableofcontents
\newpage

% --- Capitoli ---
\section{Fondamenti: Il Modello Formale (Lez. 01-02)}

\subsection{Astrazione del Corso}
Il corso di Laboratorio I (Docenti: J. Soldani, G. Vinciguerra) utilizza \textbf{JavaScript} non come strumento per il web, ma come linguaggio formale per l'implementazione di algoritmi e strutture dati.
L'obiettivo è puramente \textbf{computazionale}: analizzare come un esecutore (l'interprete) elabora simboli e valori secondo regole logiche rigorose.
L'ambiente di riferimento è astratto (l'interprete Node.js), svincolato dalle logiche del browser (DOM, interfacce grafiche).

\subsection{Dominio dei Valori (Tipi Primitivi)}
In termini insiemistici, JavaScript definisce domini di valori disgiunti, detti \textit{tipi primitivi}. Un esecutore formale opera su questi insiemi:

\begin{itemize}
    \item $\mathbb{N}$ e $\mathbb{R}$ unificati nel tipo \texttt{Number} (IEEE 754). Include elementi singolari come $\infty$ (\texttt{Infinity}) e l'elemento non-numerico assorbente \texttt{NaN}.
    \item L'algebra di Boole rappresentata dal tipo \texttt{Boolean} ($\{true, false\}$).
    \item L'insieme delle sequenze di caratteri rappresentato dal tipo \texttt{String}.
\end{itemize}

Esistono inoltre due valori singolari per la gestione dell'\textit{assenza di informazione}:
\begin{itemize}
    \item \texttt{undefined}: Valore assiomatico di default per simboli dichiarati ma non legati a un valore.
    \item \texttt{null}: Elemento neutro esplicito, denota l'insieme vuoto o l'assenza di un oggetto.
\end{itemize}

\subsection{Algebra degli Operatori e Coercizione}
Le operazioni tra tipi eterogenei seguono un'algebra di riduzione definita, nota come \textbf{Coercizione}.
Data un'operazione binaria $\otimes(A, B)$, se $type(A) \neq type(B)$, l'interprete applica funzioni di trasformazione implicita $f: T_1 \to T_2$ per rendere l'operazione decidibile.

\begin{enumerate}
    \item \textbf{Assioma della Concatenazione:} Se esiste un operando $S \in \texttt{String}$ e l'operatore è $+$, l'operazione diviene concatenazione di caratteri.
    \\ $x + "y" \Rightarrow string(x) \circ "y"$.
    \item \textbf{Assioma Aritmetico:} Per operatori $\{-, *, /, \%\}$, ogni operando viene mappato nel dominio \texttt{Number}.
    \\ $"5" - 3 \Rightarrow 5 - 3 = 2$.
\end{enumerate}

\subsection{Logica Booleana Estesa}
L'operatore di uguaglianza è definito in due modalità:
\begin{itemize}
    \item \textbf{Uguaglianza Debole ($==$):} Equivalenza logica post-riduzione (ammette coercizione).
    \item \textbf{Uguaglianza Stretta ($===$):} Identità matematica. $A \equiv B \iff (val(A) = val(B)) \land (type(A) = type(B))$.
\end{itemize}

L'insieme dei valori "Veri" e "Falsi" è esteso. L'insieme dei \textbf{Falsy Values} $\mathcal{F}$ contiene esattamente sei elementi:
$$ \mathcal{F} = \{ \text{false}, 0, \text{""}, \text{null}, \text{undefined}, \text{NaN} \} $$
Ogni altro valore appartiene all'insieme complementare $\mathcal{T}$ (Truthy).
 \newpage
\section{Flusso di Controllo e Predicati (Lezione 03)}

\subsection{Sintassi Formale e Binding}
La definizione sintattica del linguaggio è descritta tramite la notazione \textbf{BNF} (Backus-Naur Form), che definisce ricorsivamente le produzioni valide.
La gestione dello stato avviene mediante il meccanismo di \textbf{binding} (legame) tra un Identificatore e un Valore in memoria:
\begin{itemize}
    \item \texttt{let}: Definisce un binding mutabile. Lo stato associato all'identificatore può variare nel tempo ($x_t \neq x_{t+1}$).
    \item \texttt{const}: Definisce un binding immutabile. L'associazione $Id \to Valore$ è costante per l'intera esecuzione.
    \item \texttt{var}: (Deprecato) Binding con scope di funzione e proprietà di \textit{hoisting}.
\end{itemize}

\subsection{Predicati e Ramificazioni}
Le strutture condizionali valutano un predicato logico $P$.
\begin{itemize}
    \item \textbf{If / Else:} Dato un predicato $P$, il flusso $\Phi$ si divide in due rami $\Phi_{true}$ e $\Phi_{false}$.
    \item \textbf{Operatore Ternario:} Funzione che mappa un booleano in due possibili valori di ritorno: $f(P) \to \{v_1, v_2\}$.
\end{itemize}

Lo \textbf{Switch} è una struttura di selezione discreta. Dato un valore $x$ e un insieme di costanti $C = \{c_1, c_2, \dots\}$, trasferisce il controllo al ramo $i$ tale che $x \equiv c_i$.
\textit{Nota:} L'assenza di un'istruzione di interruzione (\texttt{break}) causa l'esecuzione sequenziale dei rami successivi ("fall-through").

\subsection{Automatizzare con i Cicli (Loops)}
Per compiere operazioni ripetitive, usiamo le strutture iterative.
Il ciclo \textbf{For} è il "coltellino swizzero": composto da inizializzazione, condizione e incremento, ci permette di controllare perfettamente il numero di iterazioni. Esistono varianti comode come \texttt{for..in} (per scorrere le chiavi di un oggetto) e \texttt{for..of} (per i valori di un array).

Se non sappiamo a priori quante volte dovremo iterare, usiamo il \textbf{While}, che continua finché la condizione è vera. Il suo fratello \textbf{Do-While} garantisce che il blocco venga eseguito \textit{almeno una volta} prima di controllare la condizione.
All'interno dei cicli, possiamo usare \texttt{break} per uscire immediatamente o \texttt{continue} per saltare al giro successivo.

\subsection{Misurare il Tempo}
Infine, una competenza pratica: come misurare quanto è veloce il nostro codice? Utilizziamo \texttt{Date.now()} per ottenere il timestamp attuale (in millisecondi). Salvando il tempo all'inizio (\texttt{start}) e alla fine (\texttt{end}), la differenza ci darà la durata dell'esecuzione.
 \newpage
\section{Problemi Algoritmici (Lezione 04)}

\subsection{Definizione e Risoluzione Algoritmi}
In questa sezione applichiamo i costrutti del linguaggio alla risoluzione di problemi computazionali classici.

\subsubsection{1. Calcolo Iterativo (Fattoriale)}
Il problema richiede il calcolo della funzione $f(n) = n!$.
L'approccio implementativo utilizza un \textbf{accumulatore di stato} ($acc$). Ad ogni passo $i$ dell'iterazione, lo stato viene aggiornato secondo la regola $acc_{i} = acc_{i-1} \times i$, con condizione di arresto $i=1$.
Questo evidenzia il concetto di invarianza di ciclo.

\subsubsection{2. Generazione di Serie (Pattern Asterischi)}
Il problema riguarda la generazione di sequenze di lunghezza variabile dipendente da un parametro $n$.
Si utilizzano \textbf{cicli annidati} (nested loops):
\begin{itemize}
    \item Il ciclo esterno controlla la cardinalità della riga $i$.
    \item Il ciclo interno esegue $k$ operazioni, dove $k$ è funzione di $i$ ($k = f(i)$).
\end{itemize}
La complessità computazionale è quadratica $O(n^2)$.

\subsubsection{3. Filtraggio Condizionale (Stampa Selettiva)}
Data una sequenza di input $S = \{x_1, x_2, \dots, 0\}$, l'algoritmo agisce come un \textbf{filtro}.
Per ogni elemento $x_k$, viene valutato un predicato composto $P(x_k, x_{k-1})$.
$$ Output(x_k) \iff (x_k > 0 \land even(x_k)) \lor (x_k < 0 \land x_k > x_{k-1}) $$
Questo richiede il mantenimento di uno stato "memoria" per l'elemento precedente $x_{k-1}$.

\subsubsection{4. Numeri Perfetti}
Un numero si dice "perfetto" quando la somma dei suoi divisori propri (escluso se stesso) è uguale al numero stesso. Ad esempio, 6 è perfetto perché $1+2+3 = 6$.
L'esercizio chiede di trovare il più piccolo numero perfetto maggiore di un dato $N$. Qui combiniamo cicli e operatori matematici (modulo \%) per testare la divisibilità.

\subsubsection{5. Analisi e Grafico di Funzione}
Gli ultimi esercizi ci portano nel mondo dell'analisi matematica applicata.
Dapprima si richiede di analizzare una funzione $f(x)$ in un intervallo $[a, b]$, trovandone gli zeri (dove incrocia l'asse X) e i punti di minimo e massimo.
Infine, la sfida suprema: disegnare il grafico della funzione usando caratteri ASCII in console, data un'area e una risoluzione. Questo richiede di mappare coordinate matematiche continue su una griglia discreta di caratteri.
 \newpage
\section{Oggetti e Array: Strutture Dati (Lezione 05)}

\subsection{Gli Oggetti (Objects)}
In JavaScript, gli oggetti sono contenitori per valori con etichetta (proprietà).
Pensa a un oggetto come a una "cosa" della vita reale. Una macchina ha proprietà come: peso, colore, modello.

\begin{concept}{Struttura: Chiave $\to$ Valore}
Pensa a un oggetto come a una "cosa" della vita reale. Una macchina ha proprietà come: peso, colore, modello.
\end{concept}
\begin{jsbox}
const auto = {
    marca: "Fiat",
    modello: "500",
    colore: "bianco"
};
\end{jsbox}

\subsubsection{Come accedere ai dati}
Ci sono due modi per leggere una proprietà:
\begin{enumerate}
    \item \textbf{Dot Notation (Punto):} \texttt{auto.marca} // "Fiat" (Più comune)
    \item \textbf{Bracket Notation (Parentesi):} \texttt{auto["marca"]} // "Fiat" (Utile se il nome della proprietà è in una variabile)
\end{enumerate}

\subsection{Memory: Valori vs Riferimenti}
Capire come varia la memoria è fondamentale. In JavaScript ci sono due regole diverse.

\subsubsection{1. Tipi Primitivi (Valori)}
\begin{concept}{Copia per Valore}
Numeri, Stringhe, Booleani. La variabile contiene \textbf{direttamente} il dato. Se copi, duplichi il dato.
\end{concept}
\begin{jsbox}
let a = 5;
let b = 5;
let c = a; // Copio il valore 5

console.log(a == b); // TRUE (5 e' uguale a 5)
console.log(a == c); // TRUE
\end{jsbox}

\subsubsection{2. Tipi Complessi (Riferimenti)}
\begin{concept}{Copia per Riferimento (Il Telecomando)}
Oggetti, Array, Funzioni.
La variabile contiene un \textbf{indirizzo di memoria} (un puntatore), non l'oggetto vero e proprio.
Se copi la variabile, copi solo il puntatore (telecomando), non l'oggetto (TV).
\end{concept}
\begin{jsbox}
let list1 = [1, 2, 3];
let list2 = [1, 2, 3];

// Sembrano uguali, ma sono due casi diversi in memoria!
console.log(list1 == list2); // FALSE (Indirizzi diversi)

let list3 = list1; // Copio l'indirizzo (il telecomando)
console.log(list1 == list3); // TRUE (Puntano allo stesso array)
\end{jsbox}

\subsubsection{Modifica Condivisa}
\begin{concept}{Side Effect (Effetto Collaterale)}
Poiché \texttt{list3} e \texttt{list1} puntano allo stesso oggetto, se ne modifico uno, cambia anche l'altro.
\end{concept}
\begin{jsbox}
list3.push(99);
console.log(list1); // [1, 2, 3, 99] -> VEDO LA MODIFICA!
\end{jsbox}

\begin{concept}{Visualizzazione Memoria: Stack vs Heap}
\begin{itemize}
    \item \textbf{Stack:} Veloce, per numeri e boolean (Valori).
    \item \textbf{Heap:} Grande e dinamico, per Oggetti e Array (Riferimenti).
\end{itemize}
\end{concept}

\begin{center}
\begin{tikzpicture}[
    node distance=0cm,
    stacknode/.style={draw, rectangle, minimum width=2.5cm, minimum height=0.8cm, fill=white},
    heapnode/.style={draw, rectangle, rounded corners, minimum width=3cm, minimum height=1.5cm, fill=niceBlue!10, align=center},
    label/.style={font=\bfseries\small, color=gray}
]
% ... (TikZ Code remained same, simplified here for brevity but assuming user wants it kept) ...
% Stack
\node[label] at (0, 3.5) {STACK (Esecuzione)};
\node[stacknode] (var1) at (0, 2.5) {age = 25};
\node[stacknode] (var2) at (0, 1.7) {name = "Luca"};
\node[stacknode, fill=jsYellow!20] (ref) at (0, 0.9) {user = <Ref \#01>};

% Heap
\node[label] at (6, 3.5) {HEAP (Memoria Oggetti)};
\node[heapnode] (obj) at (6, 1) {
    \textbf{Object \#01} \\
    \{ \\
    name: "Marco", \\
    active: true \\
    \}
};

% Puntatore
\draw[-{Stealth[length=3mm]}, thick, color=niceBlue] (ref.east) to[out=0, in=180] node[midway, above, font=\tiny\ttfamily] {Puntatore} (obj.west);
\end{tikzpicture}
\end{center}

\subsection{Visualizzazione Array}
\begin{center}
\begin{tikzpicture}
    % Definizione stile matrice array
    \matrix (A) [matrix of nodes, nodes={draw, minimum size=1cm, anchor=center, fill=white}, column sep=-\pgflinewidth]
    { 10 & 45 & 22 & 99 & 7 \\ };
    
    % Indici
    \foreach \i [count=\xi from 0] in {1,...,5}
        \node[above=0.1cm of A-1-\i, font=\ttfamily\tiny, color=gray] {\xi};
        
    % Etichetta
    \node[left=0.5cm of A, font=\bfseries] {const arr = };
    
    % Freccia puntatore
    \draw[<-] (A-1-3.south) -- ++(0, -0.5) node[below, font=\small] {arr[2]};
\end{tikzpicture}
\end{center}

\subsection{Array (Vettori)}
\begin{concept}{Lista Ordinata}
Un array è una variabile speciale che può contenere più valori in una lista ordinata $(0, 1, 2, \dots)$.
\end{concept}
\begin{jsbox}
const frutta = ["Mela", "Banana", "Pera"];
console.log(frutta[0]); // "Mela"
\end{jsbox}

\subsubsection{Aggiungere e Rimuovere Elementi}
\begin{concept}{Stack Operations}
Immagina l'array come una pila di piatti.
\begin{itemize}
    \item \textbf{push(x):} Aggiunge in fondo. (Veloce)
    \item \textbf{pop():} Toglie dal fondo. (Veloce)
\end{itemize}
\end{concept}
\begin{jsbox}
let numeri = [1, 2];
numeri.push(3);    // [1, 2, 3]
numeri.pop();      // [1, 2]
numeri.unshift(0); // [0, 1, 2]
\end{jsbox}

\subsection{Esercizi di Applicazione}
Per consolidare questi concetti, la lezione propone due sfide:
\begin{enumerate}
    \item \textbf{Gestione Inventario:} Creare un sistema che accumula coppie "Prodotto-Quantità" inserite dall'utente, fermandosi quando si digita "stop". Qui l'oggetto funge da database in memoria.
    \item \textbf{Media Voti:} Un programma che permette di inserire voti per vari studenti e, al termine, calcola la media per ciascuno. Questo esercizio combina probabilmente array (per i voti) e oggetti (per associare voti agli studenti).
\end{enumerate}
 \newpage
\section{Funzioni, Scope e Stack (Lez. 06-07)}

\subsection{Cosa sono le Funzioni}
Una funzione è un blocco di codice progettato per eseguire un compito specifico. La esegui "chiamandola" (invocandola).
Serve a non dover riscrivere lo stesso codice tante volte.

\begin{concept}{Definizione Funzione}
Una funzione è un blocco di codice progettato per eseguire un compito specifico. La esegui "chiamandola" (invocandola).
\end{concept}
\begin{jsbox}
// Definizione
function saluta(nome) {
    return "Ciao " + nome + "!";
}

// Utilizzo
let messaggio = saluta("Mario"); 
console.log(messaggio); // Stampa: "Ciao Mario!"
\end{jsbox}

\subsection{Scope: Dove vivono le variabili}
In JavaScript, le variabili non sono visibili ovunque. Dipende da dove le dichiari.

\subsubsection{1. Block Scope (let, const)}
\begin{concept}{Scope di Blocco (Recinto Privato)}
Le variabili definite con \texttt{let} o \texttt{const} dentro un blocco \texttt{\{ \}} (come un if o un for) esistono \textbf{solo lì dentro}.
\end{concept}
\begin{jsbox}
// Scope limitato (dentro un blocco if)
if (true) {
    let x = 2;
}
// x qui non esiste piu'
\end{jsbox}

\subsubsection{2. Function Scope (var)}
\begin{concept}{Scope di Funzione (Old School)}
Le variabili definite con \texttt{var} dentro una funzione sono visibili in tutta la funzione, ma non fuori.
\end{concept}
\begin{jsbox}
function test() {
    var car = "Fiat";
}
// car qui non esiste
\end{jsbox}

\subsection{Arrow Functions}
\begin{concept}{Sintassi Concisa}
Le Arrow Functions \texttt{() =>} semplificano la scrittura. Se c'è una sola espressione, il \texttt{return} è implicito.
\end{concept}
\begin{jsbox}
// Prima
let hello = function() {
    return "Hello World!";
}

// Dopo (Arrow Function)
let hello = () => {
    return "Hello World!";
}

// Se hai una sola riga, puoi togliere le graffe e il return:
let hello = () => "Hello World!";
\end{jsbox}

\subsection{Closures (Chiusure)}
\begin{concept}{Memoria Persistente}
Una closure è quando una funzione figlia "ricorda" le variabili dello scope padre anche dopo che il padre è terminato.
\end{concept}
\begin{jsbox}
function creaContatore() {
    let count = 0;
    return function() {
        count += 1; // Ricorda 'count'
        return count;
    }
}

const conta = creaContatore();
console.log(conta()); // 1
console.log(conta()); // 2
console.log(conta()); // 3
\end{jsbox}

\subsection{Ricorsione}
\begin{concept}{Funzione che chiama se stessa}
È utile per risolvere problemi che si possono spezzare in problemi più piccoli uguali (come il Countdown).
\end{concept}
\begin{jsbox}
function countdown(n) {
    if (n <= 0) {
        console.log("Via!");
    } else {
        console.log(n);
        countdown(n - 1); // Chiama se stessa con n-1
    }
}
countdown(3);
// Output: 3, 2, 1, Via!
\end{jsbox}
 \newpage
\section{Il Sottobosco: Binario e Numeri (Rif: Slide Lezione 08)}

\section{Stringhe, JSON e Metodi Funzionali (Lez. 08-09)}

\subsection{Metodi per le Stringhe}
Le stringhe in JS sono primitive, ma hanno molti metodi utili per lavorarci.
Ricorda: le stringhe \textbf{non} cambiano mai (sono immutabili). Ogni metodo ne restituisce una nuova.

\begin{concept}{Immutabilità delle Stringhe}
Le stringhe in JS sono primitive e non cambiano mai. I metodi ritornano sempre una **nuova** stringa.
\end{concept}
\begin{jsbox}
let testo = "Ciao Mondo";

// Estrazione
let parte = testo.slice(0, 4); // "Ciao"
let sub = testo.substring(0, 4); // "Ciao"

// Ricerca
let pos = testo.indexOf("Mondo"); // 5

// Trasformazione
let maiusc = testo.toUpperCase(); // "CIAO MONDO"
let pulito = "  ciao  ".trim(); // "ciao" (toglie spazi ai lati)

// Padding (utile per formattare)
let num = "5";
console.log(num.padStart(4, "0")); // "0005"
\end{jsbox}

\subsection{JSON (JavaScript Object Notation)}
\begin{concept}{Serializzazione}
Il JSON è un formato testuale standard. Serve per convertire oggetti in stringhe (per inviarli) e viceversa.
\end{concept}
\begin{jsbox}
let studente = { nome: "Luca", eta: 20 };

// Da Oggetto a Stringa (per inviarlo o salvarlo)
let stringaJSON = JSON.stringify(studente);
console.log(stringaJSON); // '{"nome":"Luca","eta":20}'

// Da Stringa a Oggetto (quando ricevi dati)
let datiRicevuti = '{"nome":"Marco","eta":22}';
let nuovoStudente = JSON.parse(datiRicevuti);
console.log(nuovoStudente.nome); // "Marco"
\end{jsbox}

\subsection{Metodi Avanzati degli Array (Map, Filter, Reduce)}
Invece di usare sempre il ciclo \texttt{for}, JavaScript offre metodi potenti per manipolare gli array.

\subsubsection{1. Map (Trasformare)}
\begin{concept}{Mappatura 1-a-1}
Crea un \textbf{NUOVO} array applicando una funzione a ogni elemento. La lunghezza non cambia.
\end{concept}
\begin{jsbox}
const numeri = [1, 2, 3, 4];
const doppi = numeri.map(x => x * 2);
console.log(doppi); // [2, 4, 6, 8]
\end{jsbox}

\subsubsection{2. Filter (Filtrare)}
\begin{concept}{Selezione Logica}
Crea un \textbf{NUOVO} array con solo gli elementi che superano un test (restituiscono true).
\end{concept}
\begin{jsbox}
const voti = [18, 25, 15, 30, 12];
const promossi = voti.filter(voto => voto >= 18);
console.log(promossi); // [18, 25, 30]
\end{jsbox}

\subsubsection{3. Reduce (Ridurre/Accumulare)}
\begin{concept}{Aggregazione N-a-1}
Riduce tutto l'array a un singolo valore (es. somma, media) usando un accumulatore.
\end{concept}
\begin{jsbox}
const prezzi = [10, 20, 30];
// tot inizia a 0. Per ogni 'p' (prezzo), facciamo tot + p
const totale = prezzi.reduce((tot, p) => tot + p, 0);
console.log(totale); // 60
\end{jsbox}

\noindent\makebox[\linewidth]{\rule{\paperwidth}{0.4pt}}

\subsection{Bitwise Operators: Manipolare i Bit}

\subsubsection*{1. IL CONCETTO}
Possiamo manipolare i singoli bit di un numero. È una roba da bassa lega (driver, protocolli di rete, flag), ma utile per capire come ragiona la macchina.

\subsubsection*{2. GLI OPERATORI}
\begin{itemize}
    \item \texttt{\&} (AND): Entrambi 1.
    \item \texttt{|} (OR): Almeno uno 1.
    \item \texttt{\^} (XOR): Diversi (uno 1 e l'altro 0).
    \item \texttt{\textasciitilde} (NOT): Inverte tutto.
\end{itemize}

\subsubsection*{3. CODICE E ANALISI}
\begin{concept}{Maschere di Bit (Flags)}
Usiamo i bit come interruttori on/off compatti.
\end{concept}
\begin{jsbox}
// Maschere di bit (Flags)
const LETTURA = 1;  // 001
const SCRITTURA = 2; // 010
const ESECUZIONE = 4; // 100

let permessi = LETTURA | SCRITTURA; // 011 (3)

// Controllo: Ho il permesso di lettura?
if (permessi & LETTURA) {
    console.log("Puoi leggere");
}
\end{jsbox}

\subsubsection*{5. TRAPPOLE E ERRORI DI DISTRAZIONE}
\begin{itemize}
    \item Scambiare \texttt{\&} (Bitwise) con \texttt{\&\&} (Logico).
    \begin{itemize}
        \item \texttt{1 \& 2} $\rightarrow$ \texttt{0} (01 \& 10 non hanno bit comuni).
        \item \texttt{1 \&\& 2} $\rightarrow$ \texttt{2} (True \&\& True $\rightarrow$ restituisce l'ultimo valore).
    \end{itemize}
\end{itemize}
 \newpage
\section{Implementazione di Insiemi e Multi-Insiemi (Lez. 09)}

\subsection{Usare gli Oggetti come Insiemi (Sets)}
In JavaScript, possiamo usare gli Oggetti per simulare gli insiemi matematici.
L'idea è semplice: usiamo le chiavi dell'oggetto per rappresentare gli elementi. Il valore non importa (spesso mettiamo 1 o true).
Questo garantisce che \textbf{non ci siano duplicati} (una chiave in un oggetto è unica).

\begin{concept}{Unicità delle Chiavi}
L'idea è semplice: usiamo le chiavi dell'oggetto per rappresentare gli elementi. Il valore non importa (spesso mettiamo 1 o true).
Questo garantisce che \textbf{non ci siano duplicati} (una chiave in un oggetto è unica).
\end{concept}
\begin{jsbox}
let frutta = {
    "mela": 1,
    "pera": 1,
    "banana": 1
};

// Verifica Appartenenza (in)
if ("mela" in frutta) {
    console.log("C'e' la mela!");
}

// Aggiungere (basta assegnare)
frutta["kiwi"] = 1;

// Rimuovere (delete)
delete frutta["pera"];
\end{jsbox}

\subsection{Operazioni sugli Insiemi: Unione e Intersezione}
Possiamo scrivere semplici funzioni per combinare due insiemi.

\subsubsection{Unione (+)}
\begin{concept}{Somma Logica (OR)}
Prendiamo due oggetti e ne creiamo un terzo che ha le chiavi di entrambi.
\end{concept}
\begin{jsbox}
function unione(A, B) {
    let C = {};
    // Copio tutte le chiavi di A
    for (let k in A) C[k] = 1;
    // Copio tutte le chiavi di B (se ci sono già, sovrascrivo, non fa nulla)
    for (let k in B) C[k] = 1;
    return C;
}
\end{jsbox}

\subsubsection{Intersezione (Solo i comuni)}
\begin{concept}{Prodotto Logico (AND)}
Tengo solo gli elementi che esistono in entrambi gli insiemi.
\end{concept}
\begin{jsbox}
function intersezione(A, B) {
    let C = {};
    for (let k in A) {
        // Se c'è anche in B, lo tengo
        if (k in B) {
            C[k] = 1;
        }
    }
    return C;
}
\end{jsbox}

\subsection{Multi-Insiemi (Contatori)}
A volte vogliamo contare \textit{quante volte} appare un elemento (es. carrello della spesa: 3 mele).
In questo caso, il valore associato alla chiave non sarà più 1, ma il numero di occorrenze.

\begin{jsbox}
let carrello = {
    "mela": 3,  // 3 mele
    "latte": 2, // 2 bottiglie
    "pane": 1
};
\end{jsbox}

\subsubsection{Da Array a Multi-Insieme}
\begin{concept}{Istogramma di Frequenza}
Un esercizio classico: dato un array di parole, contare quante volte appare ciascuna.
\end{concept}
\begin{jsbox}
function contaParole(lista) {
    let contatore = {};
    for (let parola of lista) {
        if (parola in contatore) {
            contatore[parola] += 1; // Incremento
        } else {
            contatore[parola] = 1; // Inizializzo
        }
    }
    return contatore;
}

let input = ["mela", "pera", "mela"];
console.log(contaParole(input)); // { mela: 2, pera: 1 }
\end{jsbox}
 \newpage
\section{Composizione di Funzioni e Callback (Lez. 10)}

\subsection{Comporre le Funzioni}
In programmazione, spesso si usa l'output di una funzione come input per un'altra. Questo si chiama \textbf{composizione}.

\begin{concept}{Matrioska di Funzioni}
In programmazione, spesso si usa l'output di una funzione come input per un'altra. Questo si chiama \textbf{composizione}.
$f(g(x))$ in matematica.
\end{concept}
\begin{jsbox}
function raddoppia(x) { return x * 2; }
function quadrato(x) { return x * x; }

// Composizione manuale: prima il quadrato, poi raddoppio
let risultato = raddoppia(quadrato(3)); 
// Passo 1: 3 * 3 = 9
// Passo 2: 9 * 2 = 18
console.log(risultato); // 18
\end{jsbox}

\subsection{Funzioni che ritornano Funzioni (Function Factory)}
\begin{concept}{Fabbrica di Funzioni}
Possiamo scrivere funzioni che \textit{creano} nuove funzioni personalizzate, "congelando" alcuni parametri (Closure).
\end{concept}
\begin{jsbox}
function creaPolinomio(a, b, c) {
    // Ritorna una NUOVA funzione che accetta solo x
    return function(x) {
        return (a * x * x) + (b * x) + c;
    }
}

// Creo una funzione specifica: 2x^2 + x + 0
let miaEquazione = creaPolinomio(2, 1, 0);

// Ora la uso solo con x
console.log(miaEquazione(3)); // (2*9) + 3 + 0 = 21
console.log(miaEquazione(10)); // (2*100) + 10 + 0 = 210
\end{jsbox}
 \newpage
\input{09_IdentitaOggetti_B} \newpage
\section{Ricorsione e Algoritmi (Lab 11)}

\subsection{La Ricorsione: Funzioni che chiamano se stesse}
Una funzione è ricorsiva quando "usa se stessa" per risolvere un problema.
Il segreto è avere sempre un **Caso Base** (dove mi fermo) e un **Passo Ricorsivo** (dove vado avanti).

\subsubsection{Esempio: Numeri Triangolari}
La somma dei numeri da 1 a N. `triangle(3) = 3 + 2 + 1 = 6`.

\begin{lstlisting}[language=JavaScript]
function triangle(n) {
    if (n == 1) return 1;          // Caso base: mi fermo a 1
    return n + triangle(n - 1);    // Passo ricorsivo: n + il resto
}

console.log(triangle(6)); // 21
\end{lstlisting}

\subsection{Lo Stack (La Pila dei Piatti)}
\begin{concept}{Visualizzazione Stack Tracing}
Ogni volta che chiami una funzione, il computer aggiunge un "piatto" (Record) sulla pila della memoria.
\begin{enumerate}
    \item Chiamo `triangle(3)` -> Stack: [3]
    \item Lui chiama `triangle(2)` -> Stack: [3, 2]
    \item Lui chiama `triangle(1)` -> Stack: [3, 2, 1] -> Ritorna 1 (POP!)
    \item `triangle(2)` riceve 1, fa 2+1=3 -> Ritorna 3 (POP!)
    \item `triangle(3)` riceve 3, fa 3+3=6 -> Ritorna 6 (POP!)
\end{enumerate}
\end{concept}

\subsection{Tail Recursion (Ricorsione in Coda)}
\begin{concept}{Ottimizzazione Coda}
Per non riempire la memoria (Stack Overflow), passo il risultato parziale come argomento.
In questo modo il computer riutilizza lo stesso spazio.
\end{concept}

\begin{jsbox}
function triangleCoda(n, s=0) {
    if (n < 1) return s;           // Ritorno l'accumulatore
    return triangleCoda(n - 1, s + n); // Passo tutto al prossimo
}
\end{jsbox}

\subsection{Esercizi Classici}

\subsubsection{1. Cifre di un Numero}
\begin{concept}{Logica: Modulo e Divisione}
Come scomporre un numero (es. 123)?
\begin{itemize}
    \item **Ultima cifra:** `n % 10`
    \item **Tolgo l'ultima:** `Math.floor(n / 10)`
\end{itemize}
\end{concept}

\textbf{A. Stampa Cifre (stampaCifreRec)}
\begin{jsbox}
function stampaCifreRec(n) {
    if (n < 10) {
        console.log(n); // Caso base
        return;
    }
    stampaCifreRec(Math.floor(n / 10)); // Stampa le precedenti
    console.log(n % 10); // Stampa l'ultima
}
\end{jsbox}

\textbf{B. Array di Cifre (arrayCifreRec)}
\begin{lstlisting}[language=JavaScript]
function arrayCifreRec(n) {
    if (n < 10) return [n];
    
    let resto = arrayCifreRec(Math.floor(n / 10));
    resto.push(n % 10);
    return resto;
}
\end{lstlisting}

\subsubsection{2. Fibonacci (N-esimo)}
Calcolare solo un numero specifico della sequenza: $F_n = F_{n-1} + F_{n-2}$.
\begin{lstlisting}[language=JavaScript]
function fibonacci(n) {
    if (n <= 1) return n; // 0 o 1
    return fibonacci(n - 1) + fibonacci(n - 2);
}
\end{lstlisting}

\subsubsection{3. Fibonacci (Array)}
Restituire l'array dei primi N numeri di Fibonacci.
\begin{lstlisting}[language=JavaScript]
function fibonacciPrimiN(n) {
    if (n == 0) return [];
    if (n == 1) return [0];
    if (n == 2) return [0, 1];
    
    // Ottengo i precedenti
    let prec = fibonacciPrimiN(n - 1);
    // Calcolo il nuovo: somma degli ultimi due
    let nuovo = prec[prec.length - 1] + prec[prec.length - 2];
    prec.push(nuovo);
    return prec;
}
\end{lstlisting}

\subsubsection{2. Ricerca Binaria (Su array ordinato)}
Invece di cercare uno per uno, dividiamo sempre a metà.
\begin{lstlisting}[language=JavaScript]
function binarySearch(arr, el, sx, dx) {
    if (sx > dx) return -1; // Non trovato
    
    let centro = Math.floor((sx + dx) / 2);
    
    if (arr[centro] == el) return centro; // Trovato!
    
    if (arr[centro] > el) {
        // Cerca nella metà sinistra
        return binarySearch(arr, el, sx, centro - 1);
    } else {
        // Cerca nella metà destra
        return binarySearch(arr, el, centro + 1, dx);
    }
}
\end{lstlisting}
 \newpage
\section{Liste Concatenate: Concetti Base (Lab 12)}

\begin{concept}{Stack (LIFO) vs Queue (FIFO)}
\textbf{Stack:} Pila di piatti. Ultimo dentro, primo fuori (LIFO).
\textbf{Queue:} Fila alle poste. Primo dentro, primo fuori (FIFO).
\end{concept}

\begin{center}
\begin{tikzpicture}[
    stackbox/.style={draw, thick, fill=niceBlue!10, minimum width=2.5cm, minimum height=0.8cm, font=\ttfamily},
    arrow/.style={-{Stealth[length=3mm]}, thick, color=codeBlue}
]
% Contenitore Stack
\draw[thick, open on top] (-1.5, 0) -- (-1.5, -3.5) -- (1.5, -3.5) -- (1.5, 0);
% Elementi
\node[stackbox] (item1) at (0, -3.0) {index 0: 10};
\node[stackbox] (item2) at (0, -2.1) {index 1: 20};
\node[stackbox, fill=jsYellow!20] (item3) at (0, -1.2) {index 2: 30};
% Frecce
\draw[arrow] (2.5, -0.5) to[out=180, in=0] node[midway, above, font=\small\bfseries] {PUSH} (item3.east);
\end{tikzpicture}
\end{center}

\subsection{Cosa sono le Liste Concatenate?}
A differenza degli Array, dove gli elementi sono vicini in memoria e usiamo gli indici (`a[0]`), nelle Liste gli elementi sono sparsi.
Ogni elemento (**Nodo**) contiene:
\begin{itemize}
    \item **val**: Il dato (numero, stringa, etc.)
    \item **next**: Un riferimento (freccia) al prossimo nodo. L'ultimo punta a `null`.
\end{itemize}

Struttura: `Testa -> [5 | ->] -> [3 | ->] -> [6 | null]`

\subsubsection{Immaginiamolo visivamente}
\begin{verbatim}
       HEAD
        |
        v
    +------+      +------+      +------+
    | 5    | ---> | 3    | ---> | 6    | ---> null
    +------+      +------+      +------+
\end{verbatim}

\subsection{Definizione in JavaScript}
\begin{concept}{Struttura Nodo}
Non serve una classe speciale, bastano gli oggetti!
Ogni nodo ha un valore e un puntatore al prossimo.
\end{concept}
\begin{jsbox}
let nodo3 = { val: 6, next: null };
let nodo2 = { val: 3, next: nodo3 };
let head  = { val: 5, next: nodo2 };
\end{jsbox}

\subsection{Operazioni Base (Ricorsive)}

\subsubsection{1. Trovare un elemento (listFind)}
\begin{concept}{Scansione Lineare}
Scorriamo la lista finché non troviamo il valore o finiamo (null).
\end{concept}
\begin{jsbox}
function listFind(nodo, valore) {
    if (!nodo) return null; // Lista finita, non trovato
    if (nodo.val === valore) return nodo; // Trovato!
    
    // Cerca nel resto della lista
    return listFind(nodo.next, valore);
}
\end{jsbox}

\subsubsection{2. Aggiungere in fondo (listPush)}
\begin{concept}{Inserimento in Coda}
Dobbiamo "scendere" fino all'ultimo nodo (quello con next: null) e attaccarne uno nuovo.
\end{concept}
\begin{jsbox}
function listPush(nodo, valore) {
    // Caso base: se il nodo è null, creo il nuovo nodo qui
    if (!nodo) {
        return { val: valore, next: null };
    }
    
    // Passo ricorsivo: ricostruisco la lista attaccando il risultato
    nodo.next = listPush(nodo.next, valore);
    return nodo;
}
\end{jsbox}

\subsubsection{3. Rimuovere l'ultimo (listPop)}
\begin{concept}{Rimozione Ultimo}
Restituisce una coppia: `[ListaAggiornata, ValoreRimosso]`.
\end{concept}
\begin{jsbox}
function listPop(nodo) {
    if (!nodo) return [null, undefined]; // Lista vuota
    
    // Se sono l'ultimo nodo (il mio next è null)
    if (!nodo.next) {
        return [null, nodo.val]; // Mi rimuovo (ritorno null al padre)
    }
    
    // Altrimenti vado avanti
    let [nuovaCoda, valRimosso] = listPop(nodo.next);
    nodo.next = nuovaCoda; // Aggiorno il mio puntatore
    return [nodo, valRimosso];
}
\end{jsbox}
 \newpage
\section{Liste Concatenate: Esercizi Avanzati (Lab 13)}

\subsection{Manipolazione di Liste}

\subsubsection{1. Elemento N-esimo (listNth)}
\begin{concept}{Accesso Sequenziale}
Trovare l'elemento all'indice `n` riducendo `n` ad ogni passo.
\end{concept}
\begin{jsbox}
function listNth(nodo, n) {
    if (!nodo) return null;
    if (n === 0) return nodo; // Arrivati!
    return listNth(nodo.next, n - 1);
}
\end{jsbox}

\subsubsection{2. Verifica Ordinamento (listIsSorted)}
\begin{concept}{Verifica Monotona}
Controllare se ogni nodo è minore del successivo.
\end{concept}
\begin{jsbox}
function listIsSorted(nodo) {
    if (!nodo || !nodo.next) return true; // 0 o 1 elemento: OK
    
    // Se il corrente è maggiore del prossimo -> Falso
    if (nodo.val > nodo.next.val) return false;
    
    // Controllo il resto
    return listIsSorted(nodo.next);
}
\end{jsbox}

\subsubsection{3. Mappare valori (listMap)}
\begin{concept}{Mappatura Ricorsiva}
Creare una NUOVA lista applicando $f(x)$ a ogni nodo.
\end{concept}
\begin{jsbox}
function listMap(nodo, f) {
    if (!nodo) return null;
    
    return {
        val: f(nodo.val),       // Trasformo il valore
        next: listMap(nodo.next, f) // Ricorro sul resto
    };
}
\end{jsbox}

\subsubsection{4. Rimuovere pezzi (listSplice)}
\begin{concept}{Taglio Chirurgico}
Rimuovere `count` nodi a partire dalla posizione `start`.
Ci sono 3 casi: Prima del taglio, Nel taglio, Fine.
\end{concept}
\begin{jsbox}
function listSplice(nodo, start, count) {
    if (!nodo) return null; // Fine lista
    
    // Caso 1: Devo ancora arrivare al punto di taglio
    if (start > 0) {
        nodo.next = listSplice(nodo.next, start - 1, count);
        return nodo;
    }
    
    // Caso 2: Sono nel punto di taglio, devo rimuovere
    if (count > 0) {
        // Salto 'nodo' e chiamo la funzione sul prossimo
        return listSplice(nodo.next, 0, count - 1);
    }
    
    // Caso 3: Finito di tagliare (count == 0)
    return nodo;
}
\end{jsbox}
 \newpage
\section{Ripasso e Approfondimento Liste (Lab 14)}

In questo laboratorio rivediamo le operazioni fondamentali sulle liste concatenate e affrontiamo esercizi ricorsivi più complessi.

\subsection{Operazioni di Base (Ripasso)}

\subsubsection{1. Ricerca (listFind)}
\begin{concept}{Trova Primo}
Trova il primo nodo che contiene \texttt{value}.
\end{concept}
\begin{jsbox}
function listFind(head, value) {
    if (!head) return null; // Non trovato
    if (head.val === value) return head; // Trovato!
    return listFind(head.next, value); // Cerca nel resto
}
\end{jsbox}

\subsubsection{2. Inserimento Centrale (listInsert)}
\begin{concept}{Chirurgia dei Puntatori}
Inserisce un nuovo valore SUBITO DOPO il nodo \texttt{x}, aggiornando i puntatori.
\end{concept}
\begin{jsbox}
function listInsert(x, value) {
    if (!x) return;
    
    // 1. Creo il nuovo nodo che punta al vecchio 'next' di x
    let nuovo = { val: value, next: x.next };
    
    // 2. Aggiorno x per puntare al nuovo
    x.next = nuovo;
}
\end{jsbox}

\subsubsection{3. Aggiunta in Coda (listPush con Insert)}
\begin{concept}{Discesa Ricorsiva}
Scende fino all'ultimo nodo e usa \texttt{listInsert}.
\end{concept}
\begin{jsbox}
function listPush(head, value) {
    if (!head) return { val: value, next: null }; // Lista vuota
    
    // Se è l'ultimo nodo (next è null)
    if (!head.next) {
        listInsert(head, value); // Appendo qui
        return head;
    }
    
    listPush(head.next, value); // Scendo ancora
    return head;
}
\end{jsbox}

\subsection{Esercizi Ricorsivi Avanzati}

\subsubsection{1. Verifica Uguaglianza (listEquals)}
\begin{concept}{Confronto Strutturale}
Due liste sono uguali se hanno la stessa lunghezza e gli stessi valori nello stesso ordine.
\end{concept}
\begin{jsbox}
function listEquals(l1, l2) {
    // Caso Base: Entrambe finite contemporaneamente -> OK
    if (!l1 && !l2) return true;
    
    // Caso Base: Una finita e l'altra no -> DIVERSE
    if (!l1 || !l2) return false;
    
    // Controllo valore corrente E ricorsione
    return (l1.val === l2.val) && listEquals(l1.next, l2.next);
}
\end{jsbox}

\subsubsection{2. Rimozione con Taglio (listSplice)}
\begin{concept}{Taglio e Ritorno}
Rimuove \texttt{deleteCount} nodi a partire da \texttt{start}.
\end{concept}
\begin{jsbox}
function listSplice(head, start, deleteCount) {
    if (!head) return null;

    // Caso 1: Non siamo ancora arrivati al punto di taglio
    if (start > 0) {
        // Ricostruisco il nodo corrente e vado avanti
        return {
            val: head.val,
            next: listSplice(head.next, start - 1, deleteCount)
        };
    }

    // Caso 2: Siamo nel punto di taglio (start == 0)
    if (deleteCount > 0) {
        // SALTO questo nodo (non lo ritorno).
        // Chiedo alla ricorsione di rimuoverne ancora (count - 1)
        return listSplice(head.next, 0, deleteCount - 1);
    }

    // Caso 3: Finito di tagliare, ritorno il resto
    return head;
}
\end{jsbox}

\subsubsection{3. Mappa Ricorsiva (listMap)}
\begin{concept}{Trasformazione}
Crea una copia trasformata della lista.
\end{concept}
\begin{jsbox}
function listMap(head, f) {
    if (!head) return null;
    return {
        val: f(head.val),             // Trasformo valore
        next: listMap(head.next, f)   // Ricorro
    };
}
\end{jsbox}
 \newpage
\section{Funzioni di Ordine Superiore e Array (Lab 16)}

\subsection{Funzioni come Valori}
In JS le funzioni sono "Cittadini di prima classe". Significa che puoi metterle dentro una variabile, passarle come argomenti o ritornarle.
Una **Funzione di Ordine Superiore** è proprio una funzione che accetta un'altra funzione come parametro.

\subsection{Metodi degli Array (Il "Coltellino Svizzero")}
Ecco i metodi fondamentali per lavorare con gli array senza usare il ciclo for.

\subsubsection{1. forEach (Esegui per tutti)}
Esegue del codice per ogni elemento. Non ritorna nulla.
\begin{lstlisting}[language=JavaScript]
let nomi = ["Mario", "Luigi"];
nomi.forEach((nome) => console.log("Ciao " + nome));
\end{lstlisting}

\subsubsection{2. map (Trasforma)}
\begin{concept}{Visualizzazione: Map}
Crea un **nuovo array** trasformando ogni elemento uno a uno.
\begin{verbatim}
   Input:  [ 1 , 2 , 3 ]
              |   |   |
             *2  *2  *2
              |   |   |
   Output: [ 2 , 4 , 6 ]
\end{verbatim}
\end{concept}

\begin{jsbox}
let num = [1, 2, 3];
let doppi = num.map(x => x * 2); // [2, 4, 6]
\end{jsbox}

\subsubsection{3. filter (Seleziona)}
\begin{concept}{Visualizzazione: Filter}
Crea un **nuovo array** solo con gli elementi che passano il test (return true).
\begin{verbatim}
   Input:  [ 10 , 5 , 20 ]
              |   |    |
             OK   NO  OK
              |        |
   Output: [ 10 ,     20 ]
\end{verbatim}
\end{concept}

\begin{jsbox}
let num = [10, 5, 20];
let grandi = num.filter(x => x > 9); // [10, 20]
\end{jsbox}

\subsubsection{4. find e findIndex (Cerca)}
\begin{concept}{Differenza Find vs FindIndex}
\texttt{find} ti dà il \textbf{valore} del primo elemento trovato. \texttt{findIndex} ti dà la sua \textbf{posizione}.
\end{concept}
\begin{jsbox}
let num = [5, 12, 8];
let trovato = num.find(x => x > 10); // 12
let indice = num.findIndex(x => x > 10); // 1
\end{jsbox}

\subsubsection{5. every e some (Verifica)}
\begin{concept}{Verifica Booleana}
Restituiscono un Booleano (True/False).
\begin{itemize}
    \item \textbf{every:} Tutti soddisfano la condizione?
    \item \textbf{some:} Almeno uno soddisfa la condizione?
\end{itemize}
\end{concept}
\begin{lstlisting}[language=JavaScript]
let num = [2, 4, 6];
let tuttiPari = num.every(x => x % 2 == 0); // true
\end{lstlisting}

\subsubsection{6. reduce (Accumula)}
Riduce tutto l'array a un singolo valore (es. somma).
\begin{lstlisting}[language=JavaScript]
let num = [1, 2, 3];
let somma = num.reduce((acc, x) => acc + x, 0); // 6
\end{lstlisting}

\subsubsection{7. sort (Ordina)}
Ordina l'array **modificandolo** (In-Place).
Attenzione: di base ordina come stringhe! Per i numeri serve una funzione di comparazione.
\begin{lstlisting}[language=JavaScript]
let numeri = [10, 2, 5];
numeri.sort((a, b) => a - b); // [2, 5, 10] (Crescente)
numeri.sort((a, b) => b - a); // [10, 5, 2] (Decrescente)
\end{lstlisting}

\subsection{Closures (I Ricordi delle Funzioni)}
Una Closure è quando una funzione figlia "ricorda" le variabili della funzione madre, anche dopo che la madre ha finito.

\begin{lstlisting}[language=JavaScript]
function creaContatore() {
    let count = 0; // Variabile nascosta nella closure
    return function() {
        count++; // La funzione interna puo' ancora vederla!
        return count;
    }
}
const c = creaContatore();
console.log(c()); // 1
console.log(c()); // 2
\end{lstlisting}
 \newpage
\section{Manipolazione Avanzata di Array e Variadicità (Lez. 17)}

\subsection{Destructuring (Spacchettare gli Array)}
Possiamo estrarre valori da un array e metterli in variabili "spacchettandoli" direttamente con le parentesi quadre.

\begin{concept}{Visualizzazione: Destructuring}
\begin{center}
\begin{tikzpicture}[
    node distance=1cm,
    varbox/.style={draw, rectangle, fill=jsYellow!30, minimum size=0.8cm, font=\ttfamily\bfseries},
    valbox/.style={draw, rectangle, fill=white, minimum size=0.8cm, font=\ttfamily},
    arrow/.style={->, thick, color=violet}
]
    % Variabili
    \node (let) {let [};
    \node[varbox, right=0.1cm of let] (x) {x};
    \node[right=0.1cm of x] (comma) {,};
    \node[varbox, right=0.1cm of comma] (y) {y};
    \node[right=0.1cm of y] (eq) {] = [};
    
    % Valori
    \node[valbox, right=0.1cm of eq] (v1) {10};
    \node[right=0.1cm of v1] (comma2) {,};
    \node[valbox, right=0.1cm of comma2] (v2) {20};
    \node[right=0.1cm of v2] (end) {];};

    % Frecce Mappatura
    \draw[arrow] (v1.north) to[out=90, in=90] (x.north);
    \draw[arrow] (v2.south) to[out=-90, in=-90] (y.south);
\end{tikzpicture}
\end{center}
\end{concept}

\begin{jsbox}
let coordinate = [10, 20];
let [x, y] = coordinate;
console.log(x); // 10
console.log(y); // 20
\end{jsbox}

\subsection{Rest Operator (...)}
I tre puntini \texttt{...} servono a dire "tutto il resto". È utilissimo per prendere la testa e la coda di una lista.

\begin{concept}{Visualizzazione: Rest Operator}
\begin{center}
\begin{tikzpicture}[
    valbox/.style={draw, rectangle, fill=white, minimum size=0.8cm},
    headbox/.style={draw, rectangle, fill=codeGreen!30, minimum size=0.8cm},
    tailbox/.style={draw, rectangle, fill=codePink!30, minimum size=0.8cm}
]
    \node[headbox] (n1) at (0,0) {1};
    \node[tailbox] (n2) at (1,0) {2};
    \node[tailbox] (n3) at (2,0) {3};
    \node[tailbox] (n4) at (3,0) {4};
    
    \node[below=0.5cm of n1, font=\bfseries\color{codeGreen}] {primo};
    
    \draw[decorate, decoration={brace, amplitude=10pt, mirror}, thick] (0.5, -0.2) -- (3.5, -0.2);
    \node[below=0.8cm of n2, xshift=0.5cm, font=\bfseries\color{codePink}] {...altri};
\end{tikzpicture}
\end{center}
\end{concept}

\begin{jsbox}
let numeri = [1, 2, 3, 4, 5];
let [primo, ...altri] = numeri;

console.log(primo); // 1
console.log(altri); // [2, 3, 4, 5]
\end{jsbox}

\subsection{Funzioni Variadiche (Argomenti Variabili)}
Ecco un esempio di funzione che usa il Rest Operator. Usiamo il nuovo **Box Viola** per le funzioni speciali.

\begin{funcbox}[title=Funzione Math Universale]
function matematica(operatore, ...numeri) {
    // 'numeri' diventa un ARRAY con tutti gli extra
    // Esempio: matematica("+", 1, 2, 3) --> numeri = [1, 2, 3]
    
    if (operatore === "+") {
        return numeri.reduce((tot, n) => tot + n, 0);
    }
    if (operatore === "*") {
        return numeri.reduce((tot, n) => tot * n, 1);
    }
}
\end{funcbox}

\begin{jsbox}
console.log(matematica("+", 1, 2, 3)); // 6
console.log(matematica("*", 2, 2, 2, 2)); // 16
\end{jsbox}
 \newpage
\section{Alberi Binari (Lab 18)}

\subsection{Cos'è un Albero Binario?}
È una struttura gerarchica (come le cartelle del PC).
Ogni elemento si chiama **Nodo** e può avere massimo due figli:
\begin{enumerate}
    \item **Left** (Sinistro)
    \item **Right** (Destro)
\end{enumerate}

\subsubsection{Immaginiamolo visivamente}
\begin{center}
\begin{tikzpicture}[
    level distance=1.5cm,
    sibling distance=2.5cm,
    every node/.style={circle, draw, minimum size=0.8cm, fill=niceBlue!10, font=\bfseries},
    edge from parent/.style={draw, thick, -latex}
]

\node {10}
    child { node {5}
        child { node {2} }
        child { node {7} }
    }
    child { node {15}
        child { node {12} }
        child { node {20} }
    };
\end{tikzpicture}
\end{center}

\subsubsection{Rappresentazione in JS}
Non servono classi complesse, bastano oggetti.
\begin{lstlisting}[language=JavaScript]
let nodoS = { val: 5, left: null, right: null };
let nodoD = { val: 10, left: null, right: null };
let root  = { val: 7, left: nodoS, right: nodoD };
\end{lstlisting}

\subsection{Come visitare un Albero (Traversals)}
Dobbiamo passare per tutti i nodi. Ci sono vari modi.

\subsubsection{1. Visita in Profondità (DFS)}
\begin{concept}{Strategie DFS}
Andiamo giù il più possibile, poi torniamo indietro (ricorsione).
\begin{itemize}
    \item **Pre-Order:** IO $\to$ SX $\to$ DX
    \item **In-Order:** SX $\to$ IO $\to$ DX
    \item **Post-Order:** SX $\to$ DX $\to$ IO
\end{itemize}
\end{concept}

\begin{jsbox}
// Pre-Order: IO -> SX -> DX
function preOrder(nodo) {
    if (!nodo) return;
    console.log(nodo.val);
    preOrder(nodo.left);
    preOrder(nodo.right);
}

// In-Order: SX -> IO -> DX
function inOrder(nodo) {
    if (!nodo) return;
    inOrder(nodo.left);
    console.log(nodo.val);
    inOrder(nodo.right);
}

// Post-Order: SX -> DX -> IO
function postOrder(nodo) {
    if (!nodo) return;
    postOrder(nodo.left);
    postOrder(nodo.right);
    console.log(nodo.val);
}
\end{jsbox}

\subsubsection{2. Visita in Ampiezza (BFS)}
\begin{concept}{Strategia BFS (Queue)}
Visitiamo livello per livello (prima tutti i figli, poi tutti i nipoti).
Serve una **Coda (Queue)** per ricordarsi quali nodi visitare.
\end{concept}

\begin{jsbox}
function bfs(root) {
    let coda = [root]; // Inizio con la radice
    while (coda.length > 0) {
        let n = coda.shift(); // Prendo il primo
        console.log(n.val);
        
        if (n.left) coda.push(n.left);   // Aggiungo i figli in fondo
        if (n.right) coda.push(n.right);
    }
}
\end{jsbox}

\subsection{Alberi Binari di Ricerca (BST)}
Un albero speciale dove per ogni nodo:
\begin{itemize}
    \item Tutto ciò che è a \textbf{Sinistra} è \textbf{Minore}.
    \item Tutto ciò che è a \textbf{Destra} è \textbf{Maggiore}.
\end{itemize}

\subsubsection{Ricerca in un BST}
È velocissima (come la Binary Search negli array).
\begin{lstlisting}[language=JavaScript]
function search(nodo, x) {
    if (!nodo) return false;      // Non trovato
    if (nodo.val == x) return true; // Trovato!
    
    if (x < nodo.val) return search(nodo.left, x); // Cerca a sinistra
    return search(nodo.right, x); // Cerca a destra
}

\subsubsection{Inserimento in un BST}
Per inserire, scendiamo finché non troviamo un posto vuoto (`null`).
\begin{lstlisting}[language=JavaScript]
function insert(nodo, val) {
    if (!nodo) return { val: val, left: null, right: null };
    
    if (val < nodo.val) {
        nodo.left = insert(nodo.left, val);
    } else {
        nodo.right = insert(nodo.right, val);
    }
    return nodo;
}
\end{lstlisting}
\end{lstlisting}
 \newpage
\section{Grafi: Struttura (Lab 20)}

\subsection{Cos'è un Grafo?}
Un Grafo è un insieme di punti (**Nodi** o Vertici) collegati da linee (**Archi** o Edges).
\begin{itemize}
    \item **Orientato:** Le linee sono frecce (A va a B, ma non viceversa).
    \item **Non Orientato:** Le linee sono strade a doppio senso (A è collegato a B).
\end{itemize}

\begin{verbatim}
    (A) -----> (B)
     |          |
     v          v
    (C) <----- (D)
\end{verbatim}

\subsection{Come salvarlo in memoria?}
\begin{concept}{Lista di Adiacenza}
Il modo migliore in JavaScript è usare una \textbf{Mappa} dove:
Chiave = Il Nodo, Valore = Lista dei vicini (Array).
\end{concept}
\begin{jsbox}
class Graph {
    constructor() {
        this.adjList = new Map();
    }

    addNode(nodo) {
        if (!this.adjList.has(nodo)) {
            this.adjList.set(nodo, []); // Creo lista vuota
        }
    }

    addEdge(partenza, arrivo) {
        this.addNode(partenza);
        this.addNode(arrivo);
        
        // Aggiungo il collegamento
        this.adjList.get(partenza).push(arrivo);
        
        // Se fosse NON orientato, dovrei aggiungere anche il viceversa:
        // this.adjList.get(arrivo).push(partenza);
    }
}

let g = new Graph();
g.addEdge("A", "B"); // A -> B
g.addEdge("A", "C"); // A -> C
\end{jsbox}
 \newpage
\section{Grafi: Visite DFS e BFS (Lab 21)}

\subsection{Attenzione ai Cicli!}
A differenza degli alberi, nei grafi puoi girare in tondo.
È \textbf{obbligatorio} usare un insieme \texttt{visited} per ricordarsi dove siamo già stati, altrimenti il programma va in loop infinito.

\subsection{DFS (Visita in Profondità)}
\begin{concept}{Labirinto (Filo di Arianna)}
Esplora un sentiero fino alla fine, poi torna indietro (**Backtracking**).
Usa la **Ricorsione** (Stack).
\end{concept}
\begin{jsbox}
function dfs(nodo, visited = new Set()) {
    // 1. Segno come visitato
    visited.add(nodo);
    console.log(nodo);
    
    // 2. Prendo i vicini dalla Mappa (vedi Cap. 20)
    let vicini = grafo.get(nodo);
    
    // 3. Per ogni vicino non ancora visto, vado in profondità
    for (let vicino of vicini) {
        if (!visited.has(vicino)) {
            dfs(vicino, visited);
        }
    }
}
\end{jsbox}

\subsection{BFS (Visita in Ampiezza)}
\begin{concept}{Onda d'Urto}
Esplora a cerchi concentrici (tutti i vicini a distanza 1, poi distanza 2...).
Usa una **Coda** (FIFO).
\end{concept}
\begin{jsbox}
function bfs(startNode) {
    let visited = new Set();
    let coda = [startNode];
    
    visited.add(startNode); // Importante: segna subito!
    
    while (coda.length > 0) {
        let nodo = coda.shift(); // Prendo il primo
        console.log(nodo);
        
        let vicini = grafo.get(nodo);
        for (let v of vicini) {
            if (!visited.has(v)) {
                visited.add(v); // Segna
                coda.push(v);   // Aggiungi alla coda
            }
        }
    }
}
\end{jsbox}

\subsection{Applicazioni}
\begin{itemize}
    \item **Componenti Connesse:** Se lanci una DFS e non visiti tutti i nodi, significa che il grafo è spezzato in isole.
    \item **Shortest Path:** La BFS trova sempre la strada col minor numero di salti.
\end{itemize}
 \newpage
\section{L'identità segreta degli oggetti}

\subsection{Introduzione e Contenuti}
Il laboratorio si focalizza sulla comprensione profonda degli oggetti in JavaScript, analizzando la differenza tra valore e riferimento, l'uso di costruttori, metodi, prototipi e classi.

\subsection{Valori e Riferimenti}
JavaScript gestisce i dati in due modi fondamentali in base alla loro natura (base o complessa).

\begin{labdefinition}[Tipi di dato e Semantica]
    \begin{itemize}
        \item \textbf{Tipi Base:} (numeri, booleani, stringhe). Hanno una \textbf{semantica per valore}. Le variabili contengono direttamente il dato.
        \item \textbf{Tipi Complessi:} (array, oggetti, funzioni). Hanno una \textbf{semantica per riferimento}. Le variabili contengono un indirizzo di memoria (riferimento).
    \end{itemize}
\end{labdefinition}

\subsection{Meccanismi di Assegnamento e Uguaglianza}
\begin{itemize}
    \item \textbf{Per Valore:} L'assegnamento copia il valore; l'uguaglianza confronta il contenuto.
    \item \textbf{Per Riferimento:} L'assegnamento copia il riferimento; l'uguaglianza (\texttt{==}) confronta se i riferimenti puntano alla stessa area di memoria (identità), non se i contenuti sono uguali.
\end{itemize}

\begin{labjsbox}[Esempio Identità vs Uguaglianza]
var a = [1,2,3]; // Riferimento alfa
var b = [1,2,3]; // Riferimento beta
var c = a;       // Copia riferimento alfa

a == b; // false (alfa != beta)
a == c; // true  (alfa == alfa)
\end{labjsbox}

\begin{labconcept}[Allocazione di Memoria]
Ogni volta che viene valutato un letterale (es. \texttt{[1,2,3]} o \texttt{\{...\}}), viene allocata una \textbf{nuova area di memoria} con un nuovo riferimento unico.
\end{labconcept}

\section{Modifica di Riferimenti e Funzioni}
La condivisione di oggetti in memoria implica che modifiche effettuate tramite un riferimento siano visibili attraverso tutti gli altri puntatori allo stesso oggetto.

\subsection{Passaggio di argomenti}
JavaScript usa il \textbf{passaggio per valore} degli argomenti. Tuttavia:
\begin{itemize}
    \item Se passiamo un tipo base, la funzione lavora su una copia locale e non può modificare la variabile originale.
    \item Se passiamo un riferimento, la funzione riceve una copia del riferimento. Può quindi modificare le proprietà dell'oggetto originale.
\end{itemize}

\begin{labjsbox}[Modifica Oggetto Riferito]
function compleanno(p) { p.età++; }
var pippo = {nome: "Pippo", età: 35}; // Riferimento alfa
compleanno(pippo); // Passa alfa
// pippo -> {nome: "Pippo", età: 36}
\end{labjsbox}

\section{Strutture Dati Complesse: Grafi}
A differenza degli alberi, dove ogni nodo compare in un solo posto (proprietà di unicità del padre), nei \textbf{grafi} lo stesso oggetto può comparire in più campi.

\begin{labdefinition}[Identità nel Grafo]
In un grafo $G = \langle N, E \rangle$:
\begin{itemize}
    \item \textbf{Nodi:} oggetti del tipo \texttt{\{ val: x \}}.
    \item \textbf{Archi:} oggetti \texttt{\{ da: n1, a: n2 \}} dove \texttt{n1} e \texttt{n2} sono riferimenti ai nodi.
\end{itemize}
Per far sì che più archi puntino allo \textbf{stesso} nodo (identità), si usano i riferimenti alle variabili.
\end{labdefinition}

\section{Funzioni, Metodi e Costruttori}

\subsection{Sintassi abbreviata dei Metodi}
Esistono tre modi per definire funzioni come proprietà di un oggetto:
\begin{labjsbox}[Sintassi Metodi]
var o = {
    inc: x => x + 1,                    // Arrow function
    dec: function(x) { return x - 1; }, // Classica
    zero(x) { return x - x; }           // Abbreviata (Metodo)
};
\end{labjsbox}

\subsection{Oggetti Globali Pre-dichiarati}
Variabili come \texttt{console} e \texttt{Math} sono \textbf{oggetti globali}.
\begin{itemize}
    \item \texttt{typeof console} $\rightarrow$ \texttt{"object"}.
    \item \texttt{Math.max()} o \texttt{console.log()} sono proprietà di tipo funzione di questi oggetti.
\end{itemize}

\begin{labcuriosita}[Dinamicità]
È possibile aggiungere funzioni personalizzate a questi oggetti (es. \texttt{Math.fattoriale = ...}) per usarle nel resto del programma.
\end{labcuriosita}

\subsection{La parola chiave \texttt{this}}
Quando una funzione è proprietà di un oggetto e vi accede internamente, viene chiamata \textbf{metodo}. Si usa \texttt{this} per riferirsi all'oggetto corrente.

\begin{labjsbox}[Uso di this]
persona.compleanno = function() { this.età++; };
pippo.compleanno(); // 'this' punta a pippo
\end{labjsbox}

\subsection{L'operatore \texttt{new}}
Crea un nuovo oggetto vuoto, imposta \texttt{this} a tale oggetto nel costruttore e lo restituisce implicitamente.
\begin{labconcept}[Convenzione Costruttori]
Le funzioni destinate a \texttt{new} iniziano con la \textbf{lettera maiuscola}.
\end{labconcept}

\section{Il Mistero dei Prototipi}
Perché un array ha il metodo \texttt{.push()} se non è definito direttamente nell'oggetto?

\begin{labdefinition}[Meccanismo di Ricerca]
Ogni oggetto ha un \textbf{prototipo}. 
\begin{enumerate}
    \item Se cerco una chiave e non c'è nell'oggetto, JavaScript la cerca nel prototipo.
    \item Si risale la \textbf{catena dei prototipi} fino a \texttt{Object.prototype}.
    \item Se la chiave non esiste in tutta la catena, il risultato è \texttt{undefined}.
\end{enumerate}
\end{labdefinition}

Il prototipo di un oggetto si ottiene con \texttt{Object.getPrototypeOf(o)} o tramite la proprietà speciale \texttt{\_\_proto\_\_}. L'oggetto prototipo degli array contiene tutti i metodi standard: \textit{push, pop, map, filter, etc.}.

\section{Conclusione: Verso le Classi}
Il meccanismo dei prototipi e dei costruttori è la base per l'uso delle \textbf{Classi} in JavaScript, tema centrale della prossima lezione.
 \newpage
\section{L'identità segreta degli oggetti: Prototipi e Classi}

\subsection{Recap: Valori, Riferimenti e Indirezione}
In JavaScript, la distinzione tra tipi base e complessi è fondamentale:
\begin{itemize}
    \item \textbf{Tipi Base:} (booleani, numeri, stringhe) $\rightarrow$ Semantica per \textbf{valore}.
    \item \textbf{Tipi Complessi:} (oggetti, array, funzioni) $\rightarrow$ Semantica per \textbf{riferimento}.
\end{itemize}

La semantica per riferimento utilizza indirizzi di memoria come valori. L'accesso ai dati avviene tramite operazioni di \textbf{indirezione} (spesso implicite in JS).

\begin{definition}{Operazioni di Indirezione}
    \begin{itemize}
        \item \texttt{oggetto.chiave}: Va in memoria all'indirizzo dell'oggetto, cerca la chiave e restituisce il contenuto.
        \item \texttt{array[indice]}: Analogo al precedente, usando l'indice come chiave.
        \item \texttt{funzione(...)}: Va in memoria a recuperare la definizione ed esegue il codice.
    \end{itemize}
\end{definition}

\subsection{Identità vs Uguaglianza: Il caso dei Grafi}
L'uso dei riferimenti è cruciale nelle strutture dati complesse come i grafi, dove lo \textbf{stesso} oggetto nodo può essere puntato da molteplici archi.

\begin{concept}{Identità}
Dire "lo stesso oggetto" (Identità) non è come dire "un oggetto uguale" (Uguaglianza).
\begin{itemize}
    \item \textbf{Uguaglianza:} Due oggetti distinti con lo stesso contenuto (es. \texttt{\{val:1\}} e \texttt{\{val:1\}}).
    \item \textbf{Identità:} Due riferimenti che puntano alla stessa area di memoria.
\end{itemize}
\end{concept}

Per modellare un grafo $G=\langle N,E\rangle$ dove più archi insistono sullo stesso nodo, è indispensabile usare variabili che contengano i riferimenti ai nodi, e riutilizzare quelle variabili nella creazione degli archi.

\subsection{Costruttori e l'Operatore \texttt{new}}
Per garantire che gli oggetti abbiano sempre la struttura corretta (campi giusti), si usano funzioni \textbf{costruttori}.

\subsubsection{Funzionamento di \texttt{new}}
L'operatore \texttt{new}, anteposto a una chiamata di funzione, altera il comportamento standard:
\begin{enumerate}
    \item \textbf{Allocazione:} Crea un nuovo oggetto vuoto (chiamiamolo \textit{o}).
    \item \textbf{Invocazione:} Chiama la funzione costruttore impostando \texttt{this = o}.
    \item \textbf{Arricchimento:} Il codice della funzione usa \texttt{this} per aggiungere proprietà e metodi a \textit{o}.
    \item \textbf{Ritorno:} Restituisce implicitamente \textit{o}.
\end{enumerate}

\begin{jsbox}[title={Confronto Costruttori}]
// 1. Factory Function (Sconsigliato - crea oggetti, non istanze legate)
function grafo(N, E) {
    return { nodi: N, archi: E, size() { return this.nodi.length } };
}

// 2. Constructor Function con 'new' (Standard pre-ES6)
function Grafo(N, E) {
    this.nodi = N;
    this.archi = E;
    // this è l'oggetto appena creato da new
}
var g = new Grafo([], []);
\end{jsbox}

\subsection{I Prototipi}
Il meccanismo dei prototipi permette agli oggetti di ereditare funzionalità senza duplicarle in memoria.

\begin{definition}{Algoritmo di Lookup delle Proprietà}
Quando si accede a \texttt{obj.key}:
\begin{enumerate}
    \item JS cerca \texttt{key} direttamente in \texttt{obj}. Se trovata, restituisce il valore.
    \item Se non trovata, cerca nel \textbf{prototipo} di \texttt{obj}.
    \item Continua a risalire la catena dei prototipi.
    \item Se arriva a un prototipo \texttt{null} (fine catena) e non ha trovato nulla, restituisce \texttt{undefined}.
\end{enumerate}
\textit{Nota: La scrittura avviene sempre sull'oggetto stesso, mai sul prototipo.}
\end{definition}

\subsubsection{Visualizzazione della Catena}
Possiamo ispezionare il prototipo con \texttt{Object.getPrototypeOf(o)} o la proprietà deprecata \texttt{\_\_proto\_\_}.

\begin{center}
\begin{tikzpicture}[
    object/.style={rectangle, draw=paletteIndigo, fill=white, thick, minimum width=2.5cm, minimum height=1.5cm, align=center, drop shadow={opacity=0.3, shadow xshift=1mm, shadow yshift=-1mm}},
    pointer/.style={->, >=stealth, thick, color=paletteRegalia},
    node distance=1.5cm and 2.5cm
]

% Oggetto array
\node[object, fill=paletteLavender!20] (arr) {\textbf{Array} \texttt{a} \\ \texttt{[1, 5, 7]} \\ \footnotesize \texttt{\_\_proto\_\_: ?}};

% Array Prototype
\node[object, right=of arr, fill=paletteLavender!40] (arrProto) {\textbf{Array.prototype} \\ \textit{push, pop, map...} \\ \footnotesize \texttt{\_\_proto\_\_: ?}};

% Object Prototype
\node[object, right=of arrProto, fill=paletteLavender!60] (objProto) {\textbf{Object.prototype} \\ \textit{toString...} \\ \footnotesize \texttt{\_\_proto\_\_: null}};

% Null
\node[draw=none, right=of objProto, color=gray] (null) {\textbf{null}};

% Collegamenti
\draw[pointer] (arr.east) -- node[above, font=\footnotesize\itshape]{prototipo} (arrProto.west);
\draw[pointer] (arrProto.east) -- node[above, font=\footnotesize\itshape]{prototipo} (objProto.west);
\draw[pointer] (objProto.east) -- (null.west);

\end{tikzpicture}
\end{center}

\begin{warning}{Enumerabilità}
Il ciclo \texttt{for (k in o)} e \texttt{Object.keys(o)} restituiscono solo le chiavi \textbf{proprie} dell'oggetto, ignorando quelle ereditate dal prototipo.
\end{warning}

\subsection{Prototipi e Costruttori: Il Doppio Salto Carpiato}
Esiste un legame simbiotico tra una funzione costruttore e i prototipi degli oggetti che essa crea.

\begin{concept}{Proprietà \texttt{.prototype}}
Ogni funzione \texttt{F} ha una proprietà \texttt{F.prototype}. Quando chiamiamo \texttt{new F()}, l'oggetto creato riceve automaticamente \texttt{F.prototype} come suo prototipo interno (\texttt{\_\_proto\_\_}).
\end{concept}

Questo permette di definire metodi una sola volta nel prototipo del costruttore, rendendoli disponibili a tutte le istanze:
\begin{jsbox}[title={Pattern Costruttore + Prototipo}]
function Persona(n, e) {
    this.nome = n; // Stato interno (diverso per ogni istanza)
    this.età = e;
}
// Comportamento condiviso (uguale per tutte le istanze)
Persona.prototype.compleanno = function() { this.età++; };

var p = new Persona("Pippo", 35);
// p.__proto__ === Persona.prototype è TRUE
\end{jsbox}

\subsection{Le Classi (Syntactic Sugar)}
Le versioni recenti di JavaScript introducono la keyword \texttt{class}. Non è un nuovo modello a oggetti, ma una sintassi più pulita sopra il sistema prototipale esistente.

\begin{jsbox}[title={Sintassi Class}]
class Persona {
    constructor(nome, eta) {
        this.nome = nome;
        this.età = eta;
    }

    compleanno() {
        this.età++;
    }
}
\end{jsbox}

\subsubsection{Cosa accade dietro le quinte?}
Anche usando \texttt{class}:
\begin{itemize}
    \item \texttt{typeof Persona} restituisce \textbf{"function"}.
    \item I metodi definiti nella classe (es. \texttt{compleanno}) finiscono in \texttt{Persona.prototype}.
    \item Le proprietà definite nel costruttore (es. \texttt{nome}) finiscono direttamente nell'istanza oggetto.
\end{itemize}

\subsection{La "Mappa Proibita" (Struttura Profonda di JS)}
Le slide finali mostrano la struttura ricorsiva e intrecciata di JS. Ecco la logica:

\begin{itemize}
    \item \textbf{Ogni funzione è un oggetto:} Quindi \texttt{Array} (il costruttore) ha come prototipo \texttt{Function.prototype}.
    \item \textbf{Function.prototype è un oggetto:} Quindi il suo prototipo è \texttt{Object.prototype}.
    \item \textbf{Object è una funzione:} Il suo prototipo è \texttt{Function.prototype}.
\end{itemize}

\begin{center}
\begin{tikzpicture}[
    node distance=2.5cm,
    entity/.style={rectangle, draw=paletteRegalia, fill=white, thick, minimum size=1.2cm, drop shadow, rounded corners, align=center},
    proto/.style={->, >=stealth, dashed, color=gray!80, thick},
    prop/.style={->, >=stealth, color=paletteIndigo, thick}
]

\node[entity] (funcProto) {Function.prototype};
\node[entity, below=of funcProto] (objProto) {Object.prototype};
\node[entity, right=of funcProto] (ArrayConst) {Array\\(Funzione)};
\node[entity, below=of ArrayConst] (ArrProto) {Array.prototype};
\node (null) [right=of objProto] {\textbf{null}};

% Relazioni
\draw[proto] (ArrayConst) -- node[above]{\tiny \_\_proto\_\_} (funcProto);
\draw[proto] (funcProto) -- node[right]{\tiny \_\_proto\_\_} (objProto);
\draw[proto] (ArrProto) -- node[below]{\tiny \_\_proto\_\_} (objProto);
\draw[proto] (objProto) -- (null);

% Relazione prototype
\draw[prop] (ArrayConst) -- node[left]{\footnotesize .prototype} (ArrProto);

\end{tikzpicture}
\end{center}
 \newpage
\section{Classi, Ereditarietà e Generatori}

\section{Il Diabolico Dualismo: Proto vs Prototype}
Una delle fonti di confusione maggiori in JavaScript è la distinzione tra `__proto__` e `prototype`.

\begin{definition}{Distinzione Fondamentale}
    \begin{itemize}
        \item \textbf{\texttt{\_\_proto\_\_}}: È una proprietà presente in \textbf{TUTTI} gli oggetti (funzioni incluse). Rappresenta il legame alla catena dei prototipi (da chi eredito?).
        \item \textbf{\texttt{prototype}}: È una proprietà presente \textbf{SOLO} nelle funzioni. È l'oggetto che verrà assegnato come `__proto__` alle istanze create invocando quella funzione con `new`.
    \end{itemize}
\end{definition}

\begin{concept}{Relazione di Identità}
Data una funzione costruttore `F` e un oggetto `o` creato con `var o = new F()`, vale sempre:
$$ \texttt{F.prototype} \equiv \texttt{o.\_\_proto\_\_} $$
Nota bene: `F.__proto__` esiste, ma è il prototipo della funzione stessa (ovvero `Function.prototype`), non c'entra con gli oggetti creati da F.
\end{concept}

\section{Classi ed Espressioni Classe}
Le classi (introdotte in ES6) sono "zucchero sintattico" sopra il modello a prototipi esistente.

\subsection{Traduzione Dietro le Quinte}
Cosa succede realmente quando dichiariamo una classe?

\begin{jsbox}[title={Classe ES6}]
class Persona {
    constructor(n, e) {
        this.nome = n;
        this.eta = e;
    }
    compleanno() { this.eta++; }
}
\end{jsbox}

\begin{jsbox}[title={Equivalente ES5 (Pseudocodice)}]
function Persona(nome, eta) {
    this.nome = nome;
    this.eta = eta;
    // this.__proto__ = Persona.prototype (implicito nel new)
}
Persona.prototype = {
    constructor: Persona,
    compleanno: function() { this.eta++; }
};
\end{jsbox}

\subsection{Espressioni Classe}
Come le funzioni, le classi possono essere anonime e assegnate a variabili:
\begin{jsbox}
var Persona = class {
    constructor(n, e) { ... }
};
\end{jsbox}

\section{Ereditarietà Avanzata}

\subsection{Meccanismo di Implementazione}
L'ereditarietà `class Studente extends Persona` non fa altro che concatenare i prototipi:
\begin{enumerate}
    \item \textbf{Istanze:} `Studente.prototype.__proto__ === Persona.prototype`. Questo permette alle istanze di Studente di accedere ai metodi di Persona.
    \item \textbf{Statici:} `Studente.__proto__ === Persona`. Questo permette alla classe Studente di ereditare i metodi statici di Persona.
\end{enumerate}

\subsection{Overriding e `super`}
Una sottoclasse può ridefinire (override) i metodi.
\begin{itemize}
    \item \textbf{super.metodo()}: Invoca il metodo dal prototipo del genitore.
    \item \textbf{super() nel costruttore}: Obbligatorio prima di usare `this`.
\end{itemize}

\begin{algo}{Esecuzione di new Sottoclasse()}
Quando si esegue `new Studente(...)`:
\begin{enumerate}
    \item Viene creato un oggetto vuoto.
    \item La catena dei costruttori viene risalita fino alla radice (`Persona`).
    \item Viene eseguito il costruttore di `Persona` (inizializza proprietà base).
    \item Viene eseguito il costruttore di `Studente` (aggiunge proprietà specifiche).
    \item L'oggetto completo viene restituito.
\end{enumerate}
\end{algo}

\section{Membri di Classe (ES2022+)}

\subsection{Campi Pubblici, Privati e Statici}
JavaScript moderno permette definizioni dichiarative.

\begin{jsbox}[title={Esempio Completo Liceale}]
class Liceale extends Studente {
    // Campo Statico (della classe)
    static quanti = 0;

    // Campo Privato (visibile solo dentro la classe)
    #iscritto = false;

    // Campo Pubblico (inizializzato su ogni istanza)
    materie = ["Mate", "Italiano"];

    constructor(nome) {
        super(nome);
        Liceale.quanti++;
    }

    unipi() { return this.#iscritto; } // OK
    
    // Accesso esterno a #iscritto -> Errore SyntaxError
}
\end{jsbox}

\section{Getter e Setter (Accessors)}
Permettono di intercettare la lettura e scrittura di proprietà.

\begin{jsbox}
let o = {
    get x() { return 1; },      // Lettura o.x restituisce sempre 1
    set x(v) { console.log(v); } // Scrittura o.x = 5 logga 5
};
\end{jsbox}

\begin{algo}{Strategia di Ricerca Proprietà}
Quando si accede a `obj.key`:
\begin{enumerate}
    \item Se la chiave esiste nell'oggetto:
    \begin{itemize}
        \item Se è un \textbf{getter}, invoca la funzione e restituisce il risultato.
        \item Se è una proprietà base, restituisce il valore.
    \end{itemize}
    \item Se non esiste, cerca nel prototipo (ricorsivamente).
    \item Se non trovata: \texttt{undefined}.
\end{enumerate}
Similmente per la scrittura (con i \textbf{setter}). I setter vengono ereditati!
\end{algo}

\section{Generatori}
I generatori sono funzioni speciali che possono sospendere l'esecuzione e riprenderla, mantenendo lo stato locale.

\subsection{Sintassi e Funzionamento}
\begin{itemize}
    \item Dichiarazione: \texttt{function* nome()} o \texttt{*metodo()}.
    \item \textbf{\texttt{yield expr}}: Restituisce \texttt{expr} al chiamante e "congela" la funzione.
    \item \textbf{\texttt{next(val)}}: Riprende l'esecuzione. Se si passa \texttt{val}, questo diventa il risultato dell'espressione \texttt{yield} che era in attesa.
\end{itemize}

\begin{jsbox}[title={Esempio Range}]
function* range(a, b) {
    let i = a;
    while(i < b) yield i++;
}
// Uso
for (var n of range(4, 8)) { console.log(n); } // 4, 5, 6, 7
\end{jsbox}

\subsection{Metodi Avanzati dei Generatori}
Oltre a `next()`, l'oggetto generatore possiede:
\begin{itemize}
    \item \textbf{\texttt{return(v)}}: Termina forzatamente il generatore (come se ci fosse un `return v` nel codice).
    \item \textbf{\texttt{throw(e)}}: Lancia un'eccezione all'interno del generatore (nel punto dello `yield`).
\end{itemize}

\subsection{Generatori Ricorsivi (`yield*`)}
Per delegare la generazione a un altro iterabile, si usa `yield*`.
\begin{jsbox}
function* merge(g1, g2) {
    yield* g1; // Emette tutti i valori di g1
    yield* g2; // Poi tutti i valori di g2
}
\end{jsbox}

\subsection{Esempio Complesso: La Poesia}
Un generatore che estrae parole da un testo.
\begin{jsbox}
class Poesia {
    #testo;
    constructor(t) { this.#testo = t; }

    *parole() {
        var i = 0;
        while (true) {
            let f = this.#testo.indexOf(" ", i); // Cerca spazio
            if (f >= 0) {
                yield this.#testo.slice(i, f); // Restituisce parola
                i = f + 1; // Avanza indice
            } else {
                yield this.#testo.slice(i); // Ultima parola
                break; // Fine
            }
        }
    }
}
var p = new Poesia("La vispa Teresa...");
console.log([...p.parole()]); // Array di parole
\end{jsbox}

\section{Esercizi: Code (Queue)}
Implementazione completa richiesta dal laboratorio.

\subsection{Classe Base: Queue}
\begin{jsbox}
class Queue {
    #array = []; // Array privato per lo storage

    put(x) {
        // Aggiunge in fondo
        this.#array[this.#array.length] = x;
    }

    get() {
        // Rimuove e restituisce dalla testa
        let x = this.#array[0];
        this.#array = this.#array.slice(1);
        return x;
    }

    size() { return this.#array.length; }
    
    empty() { return this.#array.length === 0; }
    
    toString() {
        return `Queue with ${this.size()} elements`;
    }
}
\end{jsbox}

\subsection{Sottoclasse: LimitedQueue}
Eredita da Queue e aggiunge un controllo sulla capacità massima.

\begin{jsbox}
class LimitedQueue extends Queue {
    #max_size; // Capacità massima

    constructor(max_size) {
        super(); // Inizializza this tramite Queue
        this.#max_size = max_size;
    }

    // Metodo aggiuntivo
    full() {
        return this.size() == this.#max_size;
    }

    // Override di put
    put(x) {
        if (this.full()) {
            console.log('Queue is full!!');
        } else {
            super.put(x); // Delega alla superclasse
        }
    }

    // Override di toString con super
    toString() {
        return super.toString() + ` of ${this.#max_size}`;
    }
}
\end{jsbox}
 \newpage

\newpage
\section{Cheat Sheet: La Bibbia di JavaScript}
Riepilogo definitivo di sintassi, metodi e algoritmi.

\subsection{1. Oggetto Math (Matematica)}
Tutte le funzioni matematiche sono metodi statici dell'oggetto \texttt{Math}.

\begin{center}
\begin{tabular}{|l|p{5cm}|l|}
\hline
\textbf{Metodo} & \textbf{Descrizione} & \textbf{Esempio} \\
\hline
\texttt{Math.floor(x)} & Arrotonda per \textbf{difetto} (pavimento) & \texttt{floor(5.9) -> 5} \\
\hline
\texttt{Math.ceil(x)} & Arrotonda per \textbf{eccesso} (soffitto) & \texttt{ceil(5.1) -> 6} \\
\hline
\texttt{Math.round(x)} & Arrotonda all'intero più \textbf{vicino} & \texttt{round(5.5) -> 6} \\
\hline
\texttt{Math.abs(x)} & Valore assoluto (toglie il segno) & \texttt{abs(-5) -> 5} \\
\hline
\texttt{Math.min(a,b)} & Restituisce il valore più basso & \texttt{min(1, 2) -> 1} \\
\hline
\texttt{Math.max(a,b)} & Restituisce il valore più alto & \texttt{max(1, 2) -> 2} \\
\hline
\texttt{Math.random()} & Numero casuale tra 0 e 1 (escl.) & \texttt{0.1234...} \\
\hline
\texttt{Math.sqrt(x)} & Radice quadrata & \texttt{sqrt(9) -> 3} \\
\hline
\texttt{Math.pow(b, e)} & Potenza ($b^e$) & \texttt{pow(2, 3) -> 8} \\
\hline
\texttt{Math.trunc(x)} & Rimuove la parte decimale (tronca) & \texttt{trunc(5.9) -> 5} \\
\hline
\end{tabular}
\end{center}

\subsection{2. Stringhe e Array}

\subsubsection{Stringhe (Immutabili)}
\begin{tabular}{|l|l|l|}
\hline
\texttt{length} & Proprietà lunghezza & \texttt{"Ciao".length -> 4} \\
\hline
\texttt{charAt(i)} & Carattere all'indice i & \texttt{"Ciao".charAt(0) -> "C"} \\
\hline
\texttt{includes(s)} & Cerca sottostringa & \texttt{"Ciao".includes("ia") -> true} \\
\hline
\texttt{indexOf(s)} & Indice prima occorrenza & \texttt{"Ciao".indexOf("z") -> -1} \\
\hline
\texttt{slice(a, b)} & Estrae da a fino a b & \texttt{"Ciao".slice(1,3) -> "ia"} \\
\hline
\texttt{split(sep)} & Da Stringa a Array & \texttt{"A,B".split(",") -> ["A","B"]} \\
\hline
\texttt{trim()} & Rimuove spazi lati & \texttt{" x ".trim() -> "x"} \\
\hline
\texttt{substring(a,b)} & Come slice (no negativi) & \texttt{"Ciao".substring(1,3) -> "ia"} \\
\hline
\texttt{toUpperCase()} & In MAIUSCOLO & \texttt{"ciao".toUpperCase() -> "CIAO"} \\
\hline
\texttt{toLowerCase()} & In minuscolo & \texttt{"Ciao".toLowerCase() -> "ciao"} \\
\hline
\texttt{replace(a, b)} & Sostituisce a con b & \texttt{"Ciao".replace("C", "M") -> "Miao"} \\
\hline
\end{tabular}

\subsubsection{Array: Ordinamento e Modifica (Attenzione ai Mutanti!)}
\begin{center}
\begin{tabular}{|l|p{6cm}|l|}
\hline
\textbf{Metodo} & \textbf{Descrizione} & \textbf{Mutante?} \\
\hline
\texttt{sort(fn)} & Ordina in-place. Default alfabetico! & \textbf{SÌ} \\
\hline
\texttt{splice(i, n)} & Toglie n elem da i. (Opz: aggiunge) & \textbf{SÌ} \\
\hline
\texttt{toSorted(fn)} & Copia ordinata (Nuovo Array) & NO \\
\hline
\texttt{concat(arr)} & Unisce due array (Nuovo Array) & NO \\
\hline
\texttt{push/pop} & Aggiunge/Rimuove in coda & \textbf{SÌ} \\
\hline
\texttt{shift/unshift} & Rimuove/Aggiunge in testa & \textbf{SÌ} \\
\hline
\texttt{slice(a, b)} & Copia sottosezione [a, b) & NO \\
\hline
\end{tabular}
\end{center}

\subsubsection{Array: Metodi Funzionali (Iterazione e Trasformazione)}
\begin{center}
\begin{tabular}{|l|p{5cm}|l|}
\hline
\textbf{Metodo} & \textbf{Descrizione} & \textbf{Return} \\
\hline
\texttt{forEach(fn)} & Esegue su ogni elemento & \texttt{undef} \\
\hline
\texttt{map(fn)} & Crea nuovo array trasformato & \textbf{Array} \\
\hline
\texttt{filter(fn)} & Filtra gli elementi (true) & \textbf{Array} \\
\hline
\texttt{find(fn)} & Primo elemento cercato & Elem \\
\hline
\texttt{some/every(fn)} & Almeno uno / Tutti (true) & Bool \\
\hline
\texttt{reduce(fn, acc)} & Accumula in singolo valore & Valore \\
\hline
\end{tabular}
\end{center}

\begin{jsbox}
// 1. SORT vs TOSORTED
let n = [10, 2, 5];
// sort() converte in stringhe! "10" viene prima di "2"
n.sort(); // [10, 2, 5] (Errato per numeri)
n.sort((a,b) => a - b); // [2, 5, 10] (Corretto) (Modifica n)

let original = [3, 1, 2];
let ordered = original.toSorted((a,b) => a - b); 
// original è ancora [3, 1, 2], ordered è [1, 2, 3]

// 2. SPLICE (Taglia e Cuci) vs SLICE (Copia)
let a = ["A", "B", "C", "D"];
a.splice(1, 2); // Rimuove 2 elementi partendo da indice 1
console.log(a); // ["A", "D"] (Modificato!)

// 3. CONCAT
let x = [1, 2];
let y = x.concat([3, 4]); // [1, 2, 3, 4] (x invariato)
\end{jsbox}

\subsection{2.5. Funzioni e Oggetti}
Ogni funzione è un oggetto con proprietà e metodi specifici.

\begin{center}
\begin{tabular}{|l|p{6.5cm}|}
\hline
\textbf{Proprietà/Metodo} & \textbf{Descrizione} \\
\hline
\texttt{Function.length} & Numero di argomenti attesi. \\
\hline
\texttt{Function.name} & Nome della funzione. \\
\hline
\texttt{call(this, args...)} & Invoca con \texttt{this} esplicito e argomenti. \\
\hline
\texttt{apply(this, [args])} & Come \texttt{call}, ma argomenti in array. \\
\hline
\texttt{bind(this)} & Restituisce \textbf{nuova funzione} legata. \\
\hline
\end{tabular}
\end{center}

\subsection{3. Classi e Prototipi (ES6+)}
\begin{center}
\begin{tabular}{|l|p{6cm}|}
\hline
\textbf{Concetto} & \textbf{Descrizione} \\
\hline
\texttt{\_\_proto\_\_} & Link interno di \textbf{ogni oggetto} al suo genitore. \\
\hline
\texttt{prototype} & Proprietà delle \textbf{funzioni} usata da \texttt{new}. \\
\hline
\texttt{class X} & Zucchero sintattico per funzione costruttore. \\
\hline
\texttt{extends Y} & Imposta la catena dei prototipi (\texttt{X.prototype.\_\_proto\_\_ = Y.prototype}). \\
\hline
\texttt{super} & Riferimento al genitore (obbligatorio nel costruttore). \\
\hline
\end{tabular}
\end{center}

\begin{jsbox}
class Cane extends Animale {
    static specie = "Mammifero"; // Campo Statico
    #fame = 100;                 // Campo Privato (Private)

    constructor(nome) {
        super(nome); // Chiama costruttore padre
    }

    mangia() { this.#fame -= 10; }
}
\end{jsbox}

\subsection{4. Iterazione e Strutture Dati}
\begin{center}
\begin{tabular}{|l|p{4cm}|p{4.5cm}|}
\hline
\textbf{Sintassi} & \textbf{Target} & \textbf{Esempio} \\
\hline
\texttt{for (let v of arr)} & Valori (Array/String) & \texttt{"A", "B", "C"} \\
\hline
\texttt{for (let k in obj)} & Chiavi (Oggetti) & \texttt{"nome", "eta"} \\
\hline
\texttt{Object.keys(o)} & Array Chiavi & \texttt{["k1", "k2"]} \\
\hline
\texttt{Object.values(o)} & Array Valori & \texttt{[10, 20]} \\
\hline
\texttt{Object.entries(o)} & Array Coppie & \texttt{[["k1",10], ["k2",20]]} \\
\hline
\end{tabular}
\end{center}

\subsection{5. Stack e Alberi (Algoritmi)}
\textbf{Stack (LIFO)}: Usa \texttt{push()} per inserire e \texttt{pop()} per estrarre.

\textbf{Alberi Binari (Visite)}:
\begin{jsbox}
// Nodo Albero
function Nodo(v) { this.val = v; this.sx = null; this.dx = null; }

// DFS (Ricorsiva: Pre-Order)
function dfs(n) {
    if(!n) return;
    console.log(n.val); // VISITA (Sposta per In/Post-Order)
    dfs(n.sx);
    dfs(n.dx);
}

// BFS (Iterativa con Coda)
function bfs(root) {
    let coda = [root];
    while(coda.length) {
        let n = coda.shift(); // Estrai Testa
        if(n) {
            console.log(n.val);
            if(n.sx) coda.push(n.sx); // Accoda Figli
            if(n.dx) coda.push(n.dx);
        }
    }
}
\end{jsbox}

\subsection{6. Snippets Salva-Vita (Pattern Comuni)}
\begin{jsbox}
// 1. COPIA E ORDINA (Senza modificare l'originale)
// Spread operator [...] crea una copia, poi sort() ordina quella.
let copiaOrdinata = [...originale].sort((a,b) => a - b);

// 2. RIMUOVI DUPLICATI
let unici = [...new Set([1, 2, 2, 1])]; // [1, 2]

// 3. DEEP COPY (Clonare Oggetti Semplici)
let clone = JSON.parse(JSON.stringify(oggetto));

// 4. CICLO AL CONTRARIO (Utile per rimuovere elementi)
for (let i = arr.length - 1; i >= 0; i--) {
   if (arr[i] < 0) arr.splice(i, 1);
}

// 5. FREQUENCY MAP (Conteggio Elementi)
const freq = {};
for (const x of arr) { freq[x] = (freq[x] || 0) + 1; }

// 6. MATRICE NxM (Safe Initialization)
// Evita che tutte le righe puntino allo stesso array!
const matrix = Array.from({ length: N }, () => Array(M).fill(0));

// 7. MEMOIZATION (Caching Recursion)
const memo = {};
function fib(n) {
    if (n in memo) return memo[n];
    if (n <= 1) return n;
    return memo[n] = fib(n-1) + fib(n-2);
}

// 8. DEEP CLONE (Moderno)
const copy = structuredClone(original);

// 9. FLATTEN (Appiattire Array)
const flat = [1, [2, [3]]].flat(Infinity); // [1, 2, 3]
\end{jsbox}



\end{document}
